\section{Proof of \Cref{sec:time-sampling-score}}
\label{sec:proof-section21}

In this section, we first present the continuous-time framework in a Gaussian
setting in \Cref{sec:gaussian-setting-1}. The general outline of the proof of
\Cref{prop:approx_gaussian_kl} is presented in
\Cref{sec:conv-results-discr}. Technical lemmas are gathered in
\Cref{sec:technical-lemmas}. The proof of \Cref{thm:control_diffusion} is
presented in \Cref{sec:main-content}.

\subsection{Gaussian setting}
\label{sec:gaussian-setting-1}

In what follows we present the Gaussian setting used in the proof of
\Cref{prop:approx_gaussian_kl}. We assume that $p_0 = \mathrm{N}(0, \Sigma)$
with $\Sigma \in \mathbb{S}_d(\rset)_+$.  Let $\rmD \in \mathcal{M}_d(\rset)_+$
a diagonal positive matrix such that $\Sigma = \rmP^\top \rmD \rmP$ with $\rmP$
an orthonormal matrix. We consider the following forward dynamics
\begin{equation}
  \rmd x_t = - x_t \rmd t + \sqrt{2} \rmd w_t ,
\end{equation}
with $\mathcal{L}(x_0) = p_0$.
We also consider the backward dynamics given by
\begin{equation}
  \rmd y_t = \{ y_t + 2\nabla \log p_{T-t}(y_t)\} \rmd t + \sqrt{2} \rmd w_t ,
\end{equation}
with $\mathcal{L}(y_0) = p_\infty = \mathrm{N}(0, \Id)$. Note that since for any
$t \in \ccint{0,T}$ and $x \in \rset^d$, $\nabla \log p_t(x) = - \Sigma_t^{-1}x$
with $\Sigma_t = \exp[-2t]\Sigma + (1-\exp[-2t]) \Id$, we have that
$(y_t)_{t \in \ccint{0,T}}$ is a Gaussian process. In particular, we can compute
the mean and the covariance matrix of $y_t$ in a closed form for any
$t \in \ccint{0,T}$. The results of \Cref{sec:gaussian-setting-continuous} will
not be used to prove \Cref{prop:approx_gaussian_kl}. However, they provide some
insights regarding the evolution of the mean and covariance of the backward
process.

\begin{proposition}
  \label{sec:gaussian-setting-continuous}
  For any $t \in \ccint{0,T}$, we have that
  $\mathcal{L}(y_t) = \mathrm{N}(0, \bar{\Sigma}_t)$ with
  \begin{equation}
    \bar{\Sigma}_t = \rmP^\top ((1 - \exp[-2t])\bar{\rmD}_t + \exp[-2t] \bar{\rmD}_t^2) \rmP,
  \end{equation}
  and
  \begin{equation}
    \bar{\rmD}_t = (\Id + (\rmD - \Id)\exp[-2(T-t)]) \oslash (\Id + (\rmD - \Id)\exp[-2T]) \eqsp . 
  \end{equation}
\end{proposition}

Note that $\bar{\rmD}_0 = \Id$ and
$\bar{\rmD}_T = \rmD \oslash (\Id + (\rmD - \Id)\exp[-2T]) \approx \rmD$. Hence,
we have $\bar{\Sigma}_T \approx \Sigma$ and therefore
$\mathcal{L}(y_T) \approx p_0$.

\begin{proof}
  First, note that for any $t \in \ccint{0,T}$ we have that
  \begin{equation}
    \textstyle{
      y_t = y_0 + \int_0^t ( \Id - 2\Sigma_{T-t}^{-1}) y_t + \sqrt{2} w_t = y_0 + \int_0^t ( \Id - 2\rmP^\top D_{T-t}^{-1} \rmP ) y_t + \sqrt{2} w_t ,
      }
  \end{equation}
  with $\rmD_{T-t} = \exp[-2(T-t)] \rmD + (1 - \exp[-2(T-t)]) \Id$. Denote
  $\{y_t^\rmP\}_{t \in \ccint{0,T}} = \{\rmP y_t\}_{t \in
    \ccint{0,T}}$. Using that $\rmP^\top \rmP = \Id$, we have that for any
  $t \in \ccint{0,T}$
  \begin{equation}
    \textstyle{
      y_t^\rmP =y_0^\rmP + \int_0^t ( \Id - 2 D_{T-t}^{-1}) y_t^\rmP + \sqrt{2} w_t^\rmP ,
      }
    \end{equation}
    where $\{w_t^\rmP\}_{t \in \ccint{0,T}} = \{\rmP w_t\}_{t \in
      \ccint{0,T}}$. Note that since $\rmP$ is orthonormal,
    $\{w_t^\rmP\}_{t \in \ccint{0,T}}$ is also a $d$-dimensional Brownian
    motion. We also have that $\mathcal{L}(y_0^\rmP) = \mathrm{N}(0,
    \Id)$. Hence for any $\{\{y_t^{\rmP, i}\}_{t\in \ccint{0,T}}\}_{i=1}^d$ is a
    collection of $d$ independent Gaussian processes, where for any
    $i \in \{1, \dots, d\}$ and $t \in \ccint{0,T}$,
    $y_t^{\rmP, i} = \langle y_t^{\rmP}, e_i \rangle$ and $\{e_i\}_{i=1}^d$ is
    the canonical basis of $\rset^d$. Let $i \in \{1, \dots, d\}$ and for any
    $t \in \ccint{0,T}$ denote $u_t^i = \expeLigne{y_t^{\rmP, i}}$ and
    $v_t^i = \expeLigne{(y_t^{\rmP, i})^2}$. We have that for any
    $t \in \ccint{0,T}$, $\partial_t u_t^i = (1 - 1/\rmD_t^i) u_t^i$ with
    $u_0 =0$ and $\rmD_t^i = \exp[-2t] \rmD_i + 1 - \exp[-2t]$. Hence, we get
    that for any $t \in \ccint{0,T}$, $u_t^i = 0$. Using It\^o's lemma we have
    that
    \begin{equation}
      \label{eq:edo}
      \partial v_t^i = \{ 2 - 4 / \rmD_{T-t}^i \} v_t^i + 2 ,
    \end{equation}
    with $v_0^i = 1$. Denote $\alpha^i_T = (\rmD^i - 1)\exp[-2T]$, we have that
    for any $t \in \ccint{0,T}$, $\rmD_{T-t}^i = 1 + \alpha^i_T
    \exp[2t]$. Therefore, we get that for any $t \in \ccint{0,T}$
    \begin{equation}
      2 - 4 / \rmD_{T-t}^i = -2 + 2 \times (2\alpha_T^i)\exp[2t] /(1 + \alpha_T^i \exp[2t]) = - 2 + 2 \partial_t \log(1 + \alpha_T^i \exp[2t]) .
    \end{equation}
    Hence, we have that for any $t \in \ccint{0,T}$
    \begin{equation}
      \textstyle{
        \int_0^t 2 - 4 / \rmD_{T-s}^i \rmd s = -2t + \log((1 + \alpha_T^i \exp[2t])^2/(1 + \alpha_T^i )^2) .
        }
      \end{equation}
      Hence, there exists $C_t^i \in \rmc^1(\ccint{0,T}, \rset)$ such that for
      any $t \in \ccint{0,T}$, $v_t^i = C_t^i \exp[-2t](1 + \alpha_T^i \exp[2t])^2/(1 + \alpha_T^i )^2$.
      Using \eqref{eq:edo}, we have that for any $t \in \ccint{0,T}$
      \begin{equation}
        \partial_t C_t^i = 2 \exp[2t] ((1 + \alpha_T^i \exp[2t])/(1 + \alpha_T^i))^{-2} = -(1/\alpha_T^i)(1+\alpha_T^i)^2 \partial_t (1 + \alpha_T^i \exp[2t])^{-1}.
      \end{equation}
      Hence, we have that for any $t \in \ccint{0,T}$
      \begin{equation}
        C_t^i = (1/\alpha_T^i)(1+\alpha_T^i)^2 [ (1 + \alpha_T^i)^{-1} - (1 + \alpha_T^i \exp[2t])^{-1}  ] + A ,
      \end{equation}
      with $A \geq 0$. Hence, we get that for any $t \in \ccint{0,T}$
      \begin{align}
        v_t^i &= (1/\alpha_T^i) \exp[-2t] (1 + \alpha_T^i \exp[2t]) [ (1 + \alpha_T^i \exp[2t])/(1 + \alpha_T^i) - 1 ] \\
        & \qquad \qquad + A \exp[-2t](1 + \alpha_T^i \exp[2t])^2/(1 + \alpha_T^i )^2 .
      \end{align}
      In addition, we have that $v_0^i=1$ and therefore $A=1$. Therefore, for any $t \in \ccint{0,T}$ we have
      \begin{align}
        v_t^i &= (1/\alpha_T^i) \exp[-2t] (1 + \alpha_T^i \exp[2t]) [ (1 + \alpha_T^i \exp[2t])/(1 + \alpha_T^i) - 1 ] \\
              & \qquad \qquad + \exp[-2t](1 + \alpha_T^i \exp[2t])^2/(1 + \alpha_T^i )^2 \\
         &=  (1 - \exp[-2t]) (1 + \alpha_T^i \exp[2t]) /(1 + \alpha_T^i)  + \exp[-2t](1 + \alpha_T^i \exp[2t])^2/(1 + \alpha_T^i )^2 ,      
      \end{align}
      which concludes the proof.
    \end{proof}

    \subsection{Convergence results for the discretization}
\label{sec:conv-results-discr}

In what follows, we denote
$(Y_k)_{k \in \{0, \dots, N-1\}} = (\twx_{t_k})_{k \in \{0, \dots, N-1\}}$, the
sequence given by \eqref{eq:backward_disc}. The following result gives an
expansion of the covariance matrix and the mean of $Y_N$, i.e. the output of
SGM, in the case where $p = \mathrm{N}(\mu, \Sigma)$.

\begin{theorem}
  \label{prop:approx_gaussian}
  Let $N \in \nset$, $\delta > 0$ and $T = N \delta$. Then, we have that
  $\twx_{t_N} \sim \mathrm{N}(\hat{\mu}_N, \hat{\Sigma}_N)$ with
  \begin{equation}
    \hat{\Sigma}_N = \Sigma + \exp[-4T] \hat{\Sigma}_T + \delta \hat{\mathrm{E}}_T + \delta^2 \mathrm{R}_{T, \delta} \eqsp , \qquad \hat{\mu}_N = \mu + \exp[-2T] \hat{\mu}_T + \delta \hat{e}_T + \delta^2 r_{T, \delta} \eqsp ,
  \end{equation}
  where
  $\hat{\Sigma}_T, \hat{\mathrm{E}}_T, \mathrm{R}_{T, \delta} \in \rset^{d
    \times d}$, $\hat{\mu}_T, \hat{e}_T, r_{T, \delta} \in \rset^d$ and
  $\normLigne{\mathrm{R}_{T, \delta}} + \normLigne{r_{T, \delta}} \leq R$ not
  dependent on $T \geq 0$ and $\delta >0$. We have that
  \begin{align}
  &\hat{\Sigma}_T =  - (\Sigma - \Id)(\Sigma \Sigma_T^{-1})^{2}  \eqsp , \label{eq:T_approx}\\
    &\textstyle{ \hat{\mathrm{E}}_T = \Id - (1/2)\Sigma^2 (\Sigma - \Id)^{-1} \log(\Sigma)  + \exp[-2T] \tilde{\mathrm{E}}_T\eqsp .} \label{eq:gamma_approx}
  \end{align}
  In addition, we have 
  \begin{align}
    &\hat{\mu}_T = - \Sigma_T^{-1}\Sigma \mu \eqsp ,  \\
    &\textstyle{ \hat{e}_T = \{-2 \Sigma^{-1} -  (1/4) \Sigma (\Sigma - \Id)^{-1} \log(\Sigma)\} \mu + \exp[-2T] \tilde{\mu}_T \eqsp ,  }
  \end{align}
  with $\tilde{\mathrm{E}}_T, \tilde \mu_T$ bounded and not dependent on $T$.
\end{theorem}

Before turning to the proof of \Cref{prop:approx_gaussian}, we state a few
consequences of this result.

\begin{corollary}
  \label{sec:conv-results-discr-1}
  Let $\{\twx_{t_k}\}_{k=0}^N$ the sequence defined by
  \eqref{eq:backward_disc}. We have that
  $\twx_{t_N} \sim \mathrm{N}(\mu_N, \Sigma_N)$ with
  \begin{align}
    &\Sigma_N = \Sigma + \delta \Sigma_\delta + \exp[-4T] \Sigma_T + \Sigma_{\delta, T} \eqsp , \\
    &\mu_N = \mu + \delta \mu_\delta + \exp[-2T] \mu_T + \mu_{\delta, T} \eqsp ,
  \end{align}
  with
  \begin{align}
    & \Sigma_T = - (\Sigma - \Id)\Sigma^2 \eqsp , \\
    & \Sigma_\delta = \Id - (1/2)\Sigma^2 (\Sigma - \Id)^{-1} \log(\Sigma) \eqsp , \\    
    & \mu_T = \Sigma \mu  , \\
    & \mu_\delta = \{-2 \Sigma^{-1}-  (1/4) \Sigma (\Sigma - \Id)^{-1} \log(\Sigma)\} \mu  \eqsp . 
  \end{align}
  In addition, we have
  $\lim_{\delta \to 0, T \to +\infty} \normLigne{\Sigma_{\delta, T}}/(\delta +
  \exp[-4T]) = 0$ and
  $\lim_{\delta \to 0, T \to +\infty} \normLigne{\mu_{\delta, T}}/(\delta +
  \exp[-2T]) = 0$
\end{corollary}

At first sight, it might appear surprising that $\Sigma^{-1}$ does not appear in
${\Sigma}_T$ and ${\mu}_T$. Note that in the extreme case where
$\Sigma =0$ and $\delta \to0$, i.e. we only consider the error associated with
the fact that $T \neq +\infty$, then we have no error. This is because in this
case the associated continuous-time process is an Ornstein-Uhlenbeck bridge which
has distribution $\mathrm{N}(\mu, 0)$ at time $T$.

We will use the following result.

\begin{lemma}
\label{lemma:kl_gaussian}
  Let $\pi_i = \mathrm{N}(\mu_i, \Sigma_i)$ for $i \in \{0,1\}$, with
  $\mu_0, \mu_1 \in \rset^d$ and $\Sigma_0, \Sigma_1 \in
  \mathbb{S}_d(\rset)_+$. Then, we have that 
  \begin{equation}
    \KLLigne{\pi_0}{\pi_1} = (1/2) \{ \log(\det(\Sigma_1)/\det(\Sigma_0)) - d + \mathrm{Tr}(\Sigma_1^{-1}\Sigma_0) + (\mu_1 - \mu_0)^\top \Sigma_1^{-1}(\mu_1 - \mu_0)\} .
  \end{equation}
\end{lemma}

In particular, applying \Cref{lemma:kl_gaussian} we have that for any
$\Sigma \in \mathbb{S}_d(\rset)_+$
\begin{equation}
\label{eq:kl_computation}
  \KLLigne{\mathrm{N}(0, \Sigma)}{\mathrm{N}(0, \Id)} = (1/2)\{-\log(\det(\Sigma)) + \mathrm{Tr}(\Sigma) -d \}.
\end{equation}


\begin{proposition}
    \label{sec:conv-results-discr-2}
    Let $\{\twx_{t_k}\}_{k=0}^N$ the sequence defined by
    \eqref{eq:backward_disc}. We have that
    $\twx_{t_N} \sim \mathrm{N}(\mu_N, \Sigma_N)$, with $\mu_N, \Sigma_N$ given
    by \Cref{sec:conv-results-discr-1}.
    We have that
    \begin{align}
      \KLLigne{\mathrm{N}(\mu, \Sigma)}{\mathrm{N}(\mu_N, \Sigma_N)} \leq \delta \absLigne{\trace(\Sigma^{-1}\Sigma_\delta)} + \exp[-4T] \absLigne{\trace(\Sigma^{-1}\Sigma_T)} + \exp[-4T]\mu^\top\Sigma \mu + E_{T, \delta} \eqsp ,
    \end{align}
    with $E_{T, \delta}$ a higher order term such that $\lim_{T \to +\infty, \delta \to 0} E_{T, \delta} /(\delta + \exp[-4T]) = 0$.
\end{proposition}

We now prove \Cref{prop:approx_gaussian}.

\begin{proof}
  For any $k$, denote $Y_k = \twx_{t_{N-k}}$.
  First, we recall that for any $k \in \{0, \dots, N-1\}$ and $x \in \rset^d$,
  $\nabla \log p_{T-k\gamma}(x) = -\Sigma_{T-k\gamma}^{-1} x$ where for any $t \in \ccint{0,T}$
    \begin{equation}
      \Sigma_t = (1 - \exp[-2t]) \Id + \exp[-2t] \Sigma \eqsp . 
    \end{equation}
Hence, we get that for any $k \in \{0, \dots, N-1\}$
\begin{equation}
  \label{eq:recur_recu}
  Y_{k+1} = ((1 + \gamma) \Id - 2\gamma \Sigma_{T - k \gamma}^{-1}) Y_k  + 2 \gamma \Sigma_{T - k \gamma}^{-1} M_{T-k \gamma}+ \sqrt{2 \gamma} Z_{k+1} \eqsp ,
\end{equation}
where for any $t \in \ccint{0,T}$, $M_{t} = \exp[-t] \mu$. Therefore, we get
that for any $k \in \{0, \dots, N\}$, $Y_k$ is a Gaussian random variable.  % In
% addition, using \eqref{eq:recur_recu}, we have that for any
% $k \in \{0, \dots, N-1\}$
% \begin{equation}
%   \expeLigne{Y_{k+1}} = ((1 + \gamma) \Id - 2\gamma \Sigma_{T - k \gamma}) \expeLigne{Y_k} \eqsp . 
% \end{equation}
% Since, $\expeLigne{Y_0} = 0$, we have that for any $k \in \{0, \dots, N\}$,
% $\expeLigne{Y_k} = 0$.
Using \eqref{eq:recur_recu}, we have that for any
$k \in \{0, \dots, N-1\}$
\begin{equation}
  \label{eq:recursion_cov}
  \expeLigne{\hat{Y}_{k+1} \hat{Y}_{k+1}^\top} = ((1 + \gamma) \Id - 2\gamma \Sigma_{T - k \gamma}^{-1}) \expeLigne{\hat{Y}_{k} \hat{Y}_{k}^\top} ((1 + \gamma) \Id - 2\gamma \Sigma_{T - k \gamma}^{-1}) + 2 \gamma \Id \eqsp ,
\end{equation}
where for any $k \in \{0, \dots, N\}$, $\hat{Y}_k = Y_k - \expeLigne{Y_k}$.
There exists $\rmP \in \rset^{d \times d}$ orthogonal such that
$\rmD = \rmP \Sigma \rmP^\top$ is diagonal. Note that for any
$k \in \{0, \dots, N-1\}$, we have that
$\Lambda_k = \rmP ((1+ \gamma) \Id - 2 \gamma \Sigma_{T-k\gamma}^{-1})
\rmP^\top$ is diagonal. For any $k \in \{0, \dots, N\}$, define
$\rmH_k = \rmP \expeLigne{\hat{Y}_{k} \hat{Y}_{k}^\top} \rmP^\top$. Note that
$\rmH_0 = \Id$.  Using \eqref{eq:recursion_cov}, we have that for any
$k \in \{0, \dots, N-1\}$
\begin{equation}
  \label{eq:recursion_cov_cov}
  \rmH_{k+1} = \Lambda_k^2 \rmH_k + 2 \gamma \Id \eqsp . 
\end{equation}
Hence, for any $k \in \{0, \dots, N\}$, $\rmH_k$ is diagonal. For any diagonal
matrix $C \in \rset^{d \times d}$ denote $\{c^1, \dots, c^d\}$ its diagonal
elements. Let $i \in \{1, \dots, d\}$. Using \eqref{eq:recursion_cov_cov}, we have
that for any $k \in \{0, \dots, N-1\}$
\begin{equation}
  h_{k+1}^i = (\lambda_k^i)^2 h_k^i + 2 \gamma \eqsp . 
\end{equation}
Using this result we have that for any $k \in \{0, \dots, N\}$
\begin{equation}
  \label{eq:unroll_recursion}
  \textstyle{
    h_k^i = (\prod_{\ell=0}^{k-1} \lambda_\ell^i)^2 + 2\gamma \sum_{\ell=0}^{k-1} (\prod_{j=0}^{\ell-1} \lambda_{k-1-j}^i)^2 = (\prod_{\ell=0}^{k-1} \lambda_\ell^i)^2 + 2\gamma \sum_{\ell=0}^{k-1} (\prod_{j=k-\ell}^{k-1} \lambda_{j}^i)^2 \eqsp . 
    }
\end{equation}
Let $k_1, k_2 \in \{0, \dots, N\}$ with $k_1 < k_2$. In what follows, we derive
an expansion of $I_{k_1, k_2} = \prod_{k=k_1}^{k_2} \lambda_k^i$ w.r.t.
$\gamma >0$. We have that
\begin{align}
  \label{eq:Ik1k2}
I_{k_1, k_2} &\textstyle{= \prod_{k=k_1}^{k_2} \lambda_k^i = \exp[\sum_{k=k_1}^{k_2} \log(\lambda_k^i)] = \exp[\sum_{k=k_1}^{k_2} \log(1 + \gamma a_k^i)]} \eqsp , 
\end{align}
where for any $k \in \{0, \dots, N\}$, $a_k^i = 1 - 2/d_{(N-k)\gamma}^i$, with
$d_{(N-k)\gamma}^i = 1 + \exp[-2(N-k)\gamma](d^i - 1)$.  Hence, there exist
$(b_{k, \gamma}^i)_{k \in \{0, \dots, N\}}$ bounded  such that for any $k \in \{0, \dots, N\}$ we have
\begin{equation}
  \log(1 + \gamma a_k^i) = \gamma a_k^i - (\gamma^2/2) (a_k^i)^2 + \gamma^3 b_{k, \gamma}^i \eqsp . 
\end{equation}
In addition, using \Cref{prop:maclaurin}, there exists
$C_{k_1, k_2}^\gamma \geq 0$ such that $ \gamma C_{k_1, k_2}^\gamma \leq C$ with
$C \geq 0$ not dependent on $k_2, k_2 \in \{0, \dots, N\}$, $\gamma >0$ and
\begin{equation}
  \textstyle{
    \sum_{k=k_1}^{k_2} \log(1 + \gamma a_k^i) = \int_{t_1}^{t_2^+} a^i(t) \rmd t - (\gamma/2) [\int_{t_1}^{t_2^+} a^i(t)^2 \rmd t + a^i(t_2^+) - a^i(t_1)] + C_{k_1, k_2}^\gamma \gamma^3  \eqsp ,
    }
  \end{equation}
  with $t_1 = k_1 \gamma$, $t_2^+ = (k_2 + 1) \gamma$ and for any
  $t \in \ccint{0,T}$, $a_t = 1 - 2/d_{T-t}^i$ with
  $d_{T-t}^i = 1 + \exp[-2(T-t)](d^i - 1)$.  Hence, using this result and
  \eqref{eq:Ik1k2}, we get that there exists $D_{k_1, k_2}^\gamma \geq 0$ such
  that $ \gamma D_{k_1, k_2}^\gamma \leq D$ with $D \geq 0$ not dependent on
  $k_2, k_2 \in \{0, \dots, N\}$, $\gamma >0$ and
\begin{equation}
  \textstyle{I_{k_1, k_2} = \exp[\int_{t_1}^{t_2^+} a^i(t) \rmd t ] - \exp[\int_{t_1}^{t_2^+} a^i(t)](\gamma/2)[\int_{t_1}^{t_2^+} a^i(t)^2 \rmd t + a^i(t_2^+) - a^i(t_1)] + \gamma^3 D_{k_1, k_2}^\gamma  \eqsp . } \label{eq:val_Ik1k2}
\end{equation}
Using this result, we get that there exists $E^\gamma_1 \geq 0$ such that
$\gamma E^\gamma_1 \leq E$ with $E \geq 0$ not dependent on $\gamma$ such that
\begin{equation}
  \textstyle{ (\prod_{\ell=0}^{N-1} \lambda_\ell^i)^2 =  \exp[2\int_0^T a^i(t) \rmd t] - \gamma \exp[2\int_0^T a^i(t) \rmd t][\int_{0}^{T} a^i(t)^2 \rmd t + a^i(T) - a^i(0)] + \gamma^3 E_1^\gamma \eqsp . } \label{eq:term_uno}
\end{equation}
Similarly, using \eqref{eq:val_Ik1k2}, there exist $E \geq 0$ and
$(E^\gamma_{2, \ell})_{\ell \in \{0, \dots, N\}}$ such that for any
$\ell \in \{0, \dots, N \}$, $E_{2, \ell}^\gamma \geq 0$ and
$\gamma E^\gamma_{2, \ell} \leq E$ with $E \geq 0$ not dependent on $\gamma$ and
$\ell$ such that
\begin{align}
  \textstyle{2 \gamma \sum_{\ell=0}^{N-1} (\prod_{j=N-\ell}^{N-1} \lambda_j^i)^2} &= \textstyle{ (2\gamma) \sum_{\ell=0}^{N-1}   \exp[2\int_{T-\ell \gamma}^T a^i(t) \rmd t]\}} \\
                                                                                  & \textstyle{ -2\gamma^2 \sum_{\ell=0}^{N-1} \{ \exp[2\int_{T- \ell \gamma}^T a^i(t) \rmd t][\int_{T- \ell \gamma}^{T} a^i(t)^2 \rmd t + a^i(T) - a^i(T-\ell \gamma)] \} } \\
  & + \textstyle{ \gamma^4 \sum_{\ell=0}^{N-1} E_{2, \ell}^\gamma \eqsp . }
\end{align}
Therefore, using \Cref{prop:maclaurin}, there exists $E_3^\gamma$ such that
$\gamma E_3^\gamma \leq E$ with $E \geq 0$ not dependent on $\gamma$ and
\begin{align}
  \textstyle{2 \gamma \sum_{\ell=0}^{N-1} (\prod_{j=N-\ell}^{N-1} \lambda_j^i)^2} &= \textstyle{2 \int_0^T \exp[2 \int_{T-t}^T a^i(s) \rmd s] \rmd t + \gamma (1 - \exp[2 \int_0^Ta^i(t) \rmd t]) } \\
                                                                                  &\textstyle{-2\gamma \int_0^T \{ \exp[2\int_{T-t}^T a^i(s) \rmd s][\int_{T- t}^{T} a^i(s)^2 \rmd t + a^i(T) - a^i(T-t)] \} \rmd t} \\
  &+ \gamma^3 E_3^\gamma \eqsp . \label{eq:term_duo}
\end{align}
% \Cref{prop:maclaurin} and \eqref{eq:unroll_recursion},
% \begin{align}
%   h_N^i &= \textstyle{ \exp[2\int_0^T a^i(t) \rmd t] - \gamma \exp[2\int_0^T a^i(t) \rmd t]}[\int_{0}^{T} a^i(t)^2 \rmd t + a^i(T) - a^i(0)] \\
%   & \qquad + 2\textstyle{\int_0^T \exp[2\int_{T-t}^T a^i(s) \rmd s] \rmd t + \gamma^3 E^\gamma }\\
%   & \qquad \textstyle{- 2\gamma \int_0^T \exp[2\int_{T-t}^T a^i(s) \rmd s][\int_{T-t}^{T} a^i(s)^2 \rmd s + a^i(T) - a^i(T-t)] \rmd t \eqsp . }
% \end{align}
% This can be rewritten as
Hence, combining \eqref{eq:term_uno} and \eqref{eq:term_duo} we get that 
\begin{equation}
  h_N^i = c_T^i - \gamma e_T^i + \gamma^3 E^\gamma \eqsp ,
\end{equation}
with
\begin{equation}
  \label{eq:ct_res}
  \textstyle{
    c_T^i  = \exp[2\int_0^T a^i(t) \rmd t] + 2\int_0^T \exp[2\int_{T-t}^T a^i(s) \rmd s] \rmd t \eqsp . 
    }
\end{equation}
and 
\begin{align}
  \label{eq:et_res}
  e_T^i  &= \textstyle{-\exp[2\int_0^T a^i(t) \rmd t][\int_{0}^{T} a^i(t)^2 \rmd t + a^i(T) - a^i(0)]} \textstyle{+ 1 - \exp[2 \int_0^T a^i(t) \rmd t ]} \\
    & \qquad \textstyle{ - 2 \int_0^T \exp[2\int_{T-t}^T a^i(s) \rmd s][\int_{T-t}^{T} a^i(s)^2 \rmd s + a^i(T) - a^i(T-t)] \rmd t \eqsp . 
    }
\end{align}
In what follows, we compute $c_T^i$ and $e_T^i$.
\begin{enumerate}[wide, labelwidth=!, labelindent=0pt, label=(\roman*)]
\item Using \Cref{lemma:integral_uno}  we have
  \begin{equation}
    \textstyle{ \exp[2\int_0^T a^i(t) \rmd t] = d^2 \exp[-2T]/(1 + \exp[-2T](d-1))^2 \eqsp . }
  \end{equation}
  In addition, using \Cref{lemma:integral_exp_uno} we have
  \begin{equation}
    \textstyle{ \int_0^T \exp[2\int_{T-t}^T a^i(s) \rmd s] = (d/2) (1 - \exp[-2T])/(1 + \exp[-2T](d-1)) \eqsp . }
  \end{equation}
  Combining these results and \eqref{eq:ct_res}, we get that
  \begin{align}
    c_T^i &= d + d^2 \exp[-2T](1 - (1 + \exp[-2T](d-1))^{-1})/(1 + \exp[-2T](d-1)) \\
    &= d + d^2 (d-1) \exp[-4T]/(1 + \exp[-2T](d-1))^2 \eqsp . 
  \end{align}
\item We conclude for $e_T^i$ using \Cref{prop:explicit_lambda} with
  $\lambda = d^i - 1$.
\end{enumerate}
This concludes the proof of \eqref{eq:gamma_approx}. Next, we compute the
evolution of the mean. Using \eqref{eq:recur_recu}, we have 
\begin{equation}
  \label{eq:recur_recu}
  \expeLigne{Y_{k+1}} = ((1 + \gamma) \Id - 2\gamma \Sigma_{T - k \gamma}^{-1}) \expeLigne{Y_k}  + 2 \gamma \expeLigne{\Sigma_{T - k \gamma}^{-1} M_{T-k \gamma}} \eqsp ,
\end{equation}
Note that for any $k \in \{0, \dots, N-1\}$, we have that
$\Lambda_k = \rmP ((1+ \gamma) \Id - 2 \gamma \Sigma_{T-k\gamma}^{-1})
\rmP^\top$ is diagonal. For any $k \in \{0, \dots, N\}$, define
$\rmH_k = \rmP \expeLigne{Y_{k}} \rmP^\top$. Note that
$\rmH_0 = 0$. For any $k \in \{0, \dots, N-1\}$ we have that
\begin{equation}
  \label{eq:recu_expectation}
  \rmH_{k+1} = \Lambda_k \rmH_k + 2 \gamma \rmD_{T- k\gamma}^{-1} \rmV_{T- k \gamma} \eqsp ,
\end{equation}
where for any $t \in \ccint{0,T}$, $\rmD_t = \rmP \Sigma_t \rmP^\top$ and
$\rmV_t = \rmP M_{t}$. Let $i \in \{1, \dots, d\}$. Using
\eqref{eq:recu_expectation}, we have for any $k \in \{0, \dots, N-1\}$
\begin{equation}
  \label{eq:recu_one_dim_exp}
  h_{k+1}^i = \lambda_k^i h_k^i + 2 \gamma v_{T-k \gamma}^i / d_{T-k\gamma}^i \eqsp . 
\end{equation}
In what follows, we define for any $t \in \ccint{0,T}$, $r(t)^i = v_{T-t}^i / d_{T-t}^i$
and note that for any $t \in \ccint{0,T}$
\begin{equation}
  \label{eq:r_express}
  r(t)^i = \exp[-(T-t)] / (1 + \exp[-2(T-t)](d^i - 1)) (\rmP \mu)^i \eqsp . 
\end{equation}
Using \eqref{eq:recu_one_dim_exp} and that $h_0^i = 0$, we have that for any $k \in \{0, \dots, N\}$
\begin{equation}
  \label{eq:unroll_recursion_exp}
  \textstyle{
    h_k^i =  2\gamma \sum_{\ell=0}^{k-1} r((k - \ell - 1)\gamma) \prod_{j=0}^{\ell-1} \lambda_{k-1-j}^i  =  2\gamma \sum_{\ell=0}^{k-1} r((k - \ell - 1)\gamma) \prod_{j=k-\ell}^{k-1} \lambda_{j}^i \eqsp . 
  }
\end{equation}
Using \eqref{eq:val_Ik1k2}, we get that there exists $D^\gamma \geq 0$ such that
$\gamma D^\gamma \leq D$ not dependent on $\gamma$ and
\begin{align}
  &h_N^i = \textstyle{2 \gamma \sum_{k=0}^{N-1}  r(T - (k + 1) \gamma) \exp[\int_{T - k \gamma}^{T} a^i(t) \rmd t ]} \\
        & \quad  \textstyle{- \gamma^2 \sum_{k=0}^{N-1} r(T - (k + 1) \gamma) \exp[\int_{T - k \gamma}^{T} a^i(t) \rmd t ][\int_{T - k \gamma}^{T} a^i(t)^2 \rmd t + a^i(T) - a^i(T - k \gamma)] } \\
        & \quad \textstyle{  + \sum_{k=0}^{N-1}  \gamma^4 D_{k, N}^\gamma} \\
&= \textstyle{2 \gamma \sum_{k=0}^{N-1}  r(T - (k + 1) \gamma) \exp[\int_{T - k \gamma}^{T} a^i(t) \rmd t ]} \\
        & \quad  \textstyle{- \gamma^2 \sum_{k=0}^{N-1} r(T - (k + 1) \gamma) \exp[\int_{T - k \gamma}^{T} a^i(t) \rmd t ][\int_{T - k \gamma}^{T} a^i(t)^2 \rmd t + a^i(T) - a^i(T - k \gamma)] } \\
  & \quad \textstyle{  + \gamma^3  D^\gamma \eqsp . }   
\end{align}
Using \Cref{prop:maclaurin}, we get that there exists $E^\gamma \geq 0$ such that
$\gamma E^\gamma \leq E$ not dependent on $\gamma$ and
\begin{align}
  &h_N^i = \textstyle{2 \gamma \sum_{k=0}^{N-1}  r(T - (k + 1) \gamma) \exp[\int_{T - k \gamma}^{T} a^i(t) \rmd t ]} \\
  & \quad  \textstyle{- \gamma \int_0^T r(T - t) \exp[\int_{T - t}^{T} a^i(s) \rmd s ][\int_{T - t}^{T} a^i(t)^2 \rmd t + a^i(T) - a^i(T - t)] \rmd t } \\
  & \quad \textstyle{  + \gamma^3  E^\gamma \eqsp . }     
\end{align}
In addition, for any $k \in \{0, \dots, N\}$, there exists $u_k \geq 0$ with
$u_k \leq u$ and $u \geq 0$ not dependent on $k$ and
\begin{equation}
  r(T-(k+1)\gamma) = r(T-k\gamma) -r'(T-k \gamma)\gamma + u_k \gamma^2 \eqsp . 
\end{equation}
Using this result, we get that exists $F^\gamma \geq 0$ such that
$\gamma F^\gamma \leq F$ not dependent on $\gamma$ and
\begin{align}
  &h_N^i = \textstyle{2 \gamma \sum_{k=0}^{N-1}  r(T - k \gamma) \exp[\int_{T - k \gamma}^{T} a^i(t) \rmd t ]} \\
  & \quad \textstyle{-2\gamma^2 \sum_{k=0}^{N-1} r'(T-k \gamma) \exp[\int_{T - k \gamma}^{T} a^i(t) \rmd t ] } \\
  & \quad  \textstyle{- \gamma \int_0^T r(T - t) \exp[\int_{T - t}^{T} a^i(s) \rmd s ][\int_{T - t}^{T} a^i(t)^2 \rmd t + a^i(T) - a^i(T - t)] \rmd t } \\
  & \quad \textstyle{  + \gamma^3  F^\gamma \eqsp . }     
\end{align}
Using \Cref{prop:maclaurin}, we get that there exists $G^\gamma \geq 0$ such that
$\gamma G^\gamma \leq G$ not dependent on $\gamma$ and
\begin{align}
  &h_N^i = \textstyle{2 \gamma \sum_{k=0}^{N-1}  r(T - k \gamma) \exp[\int_{T - k \gamma}^{T} a^i(t) \rmd t ]} \\
  & \quad \textstyle{-2\gamma \int_0^T r'(T-t) \exp[\int_{T - t}^{T} a^i(t) \rmd s ] \rmd t } \\
  & \quad  \textstyle{- \gamma \int_0^T r(T - t) \exp[\int_{T - t}^{T} a^i(s) \rmd s ][\int_{T - t}^{T} a^i(t)^2 \rmd t + a^i(T) - a^i(T - t)] \rmd t } \\
  & \quad \textstyle{  + \gamma^3  G^\gamma \eqsp . }     
\end{align}
In addition, using \Cref{prop:maclaurin}, we get that there exists
$H^\gamma \geq 0$ such that $\gamma H^\gamma \leq H$ not dependent on $\gamma$
and
\begin{align}
  &h_N^i = \textstyle{2 \int_0^T  r(T - t) \exp[\int_{T - t}^{T} a^i(s) \rmd s ] \rmd t} \\
  &\quad \textstyle{- \gamma \{r(0) \exp[\int_{0}^{T} a^i(t) \rmd t ]- r(T)\}}  \textstyle{-2\gamma \int_0^T r'(T-t) \exp[\int_{T - t}^{T} a^i(s) \rmd s ] \rmd t } \\
  & \quad  \textstyle{- \gamma \int_0^T r(T - t) \exp[\int_{T - t}^{T} a^i(s) \rmd s ][\int_{T - t}^{T} a^i(t)^2 \rmd t + a^i(T) - a^i(T - t)] \rmd t } \\
  & \quad \textstyle{  + \gamma^3  H^\gamma \eqsp . } \label{eq:inter_integ}    
\end{align}
In addition, we have by integration by part
\begin{align}
  &\textstyle{\int_0^T r'(T-t) \exp[\int_{T - t}^{T} a^i(s) \rmd s ] \rmd t} \\
  & \qquad \qquad \textstyle{= - \{r(0) \exp[\int_0^T a^i(t) \rmd t] - r(T)\} - \int_0^T r(T-t) a^i(T-t) \exp[\int_{T - t}^{T} a^i(s) \rmd s ] \rmd t \eqsp .  }
\end{align}
Combining this result and \eqref{eq:inter_integ} we get that
\begin{align}
  &h_N^i = \textstyle{2 \int_0^T  r(T - t) \exp[\int_{T - t}^{T} a^i(s) \rmd s ] \rmd t} \\
  &\quad \textstyle{+ \gamma \{r(0) \exp[\int_{0}^{T} a^i(t) \rmd t ]- r(T)\}}  \\
  & \quad  \textstyle{- \gamma \int_0^T r(T - t) \exp[\int_{T - t}^{T} a^i(s) \rmd s ][\int_{T - t}^{T} a^i(t)^2 \rmd t + a^i(T) - 3 a^i(T - t)] \rmd t } \\
  & \quad \textstyle{  + \gamma^3  H^\gamma \eqsp . }  \label{eq:almost_dooone}
\end{align}
In what follows, we assume that $d^i \neq 0$. The case where $d^i = 0$ is left
to the reader.  Finally using \eqref{eq:r_express} and \Cref{lemma:integral_uno}
we have that for any $t \in \ccint{0,T}$
\begin{align}
  \textstyle{
  \exp[\int_{T - t}^{T} a^i(s) \rmd s ] r(T-t)^i} &= \exp[-2t] / (1 + \exp[-2t](d^i - 1))^2 (\rmP \mu)^i d^i \\
  &= \textstyle{ \exp[2 \int_{T - t}^{T} a^i(s) \rmd s ] (\rmP \mu)^i / d^i \eqsp .
    }
\end{align}
Therefore, combining this result and \eqref{eq:almost_dooone}, we get that 
\begin{align}
  &h_N^i   = (\rmP \mu)^i/d^i [\textstyle{2 \int_0^T \exp[2 \int_{T-t}^{T} a^i(s) \rmd s] \rmd t} \\
  &\quad \textstyle{+ \gamma \{\exp[2 \int_{0}^{T} a^i(t) \rmd t] - 1\}}  \\
  & \quad  \textstyle{- \gamma \int_0^T \exp[2\int_{T-t}^T a^i(s) \rmd s] [\int_{T - t}^{T} a^i(t)^2 \rmd t + a^i(T) - 3 a^i(T - t)] \rmd t }] \\
  & \quad \textstyle{  + \gamma^3  H^\gamma \eqsp , }  \label{eq:almost_doooone}
\end{align}
which concludes the proof upon using \Cref{lemma:integral_exp_uno} and \Cref{prop:mean_mean}.
\end{proof}

\subsection{Technical lemmas}
\label{sec:technical-lemmas}

We are going to make use of the following lemma which is a direct consequence of
the Euler-MacLaurin formula.

\begin{proposition}
  \label{prop:maclaurin}
  Let $f \in \rmc^\infty(\ccint{0,T})$, and
  $(u_k^\gamma)_{k \in \{0, \dots, N-1\}}$ with $N \in \nset$ and
  $\gamma = T/N > 0$ such that for any $k \in \{0, \dots, N-1\}$,
  $u_k^\gamma = f(k \gamma)$. Then, there exists $C \geq 0$ such that
  \begin{equation}
    \textstyle{
      \int_0^T f(t) \rmd t - \gamma \sum_{k=0}^{N-1} u_k^\gamma -  (\gamma/2) \{f(T) - f(0)\} =   C \gamma^2 \eqsp .
      }
  \end{equation}
\end{proposition}

\begin{proof}
  Apply the classical Euler-MacLaurin formula to $t \mapsto f(t \gamma)$.
\end{proof}

% \begin{proof}
%   \textcolor{red}{Ref MacLaurin}
% \end{proof}

We will also use the following lemmas.

\begin{lemma}
  \label{lemma:integral_uno}
  Let $\lambda \in \ooint{-1, +\infty}$ and $a: \ \ccint{0,T} \to \rset$ such
  that for any $t \in \ccint{0,T}$,
  \begin{equation}
    a(t) = 1 - 2/(1 + \exp[-2(T-t)]\lambda) \eqsp . 
  \end{equation}
  Then, we have that for any $t \in \ccint{0,T}$,
  \begin{equation}
    \textstyle{
      \int_{T-t}^T a(s) \rmd s = t + \log((1+\lambda)/(\exp[2t]+ \lambda)) \eqsp . 
      }
    \end{equation}
    In particular, we have that for any $t \in \ccint{0,T}$
    \begin{equation}
          \textstyle{
      \exp[2\int_{T-t}^T a(s) \rmd s] = \exp[-2t] (1+\lambda)^2 /(1+ \lambda\exp[-2t])^2 \eqsp . 
      }
    \end{equation}

  \end{lemma}

  \begin{proof}
    Let $t \in \ccint{0,T}$. We have that
    $\int_{T-t}^T a(s) \rmd s = \int_0^{t} a(T-s) \rmd s$.  Define $b$ such that
    for any $t \in \ccint{0,T}$, $b(t) = a (T-t)$. In particular, we have that
    for any $t \in \ccint{0,T}$
    \begin{equation}
      b(t) = 1 - 2 / (1 +\lambda \exp[-2t]) \eqsp . 
    \end{equation}
    Hence, we have
    \begin{align}
      \textstyle{\int_0^t b(s) \rmd s} &= \textstyle{t - 2 \int_0^t (1 + \lambda \exp[-2s])^{-1} \rmd s} \\
                                       &= \textstyle{t - \int_0^t 2 \exp[2s] / (\exp[2s] + \lambda ) \rmd s} \\
                                       &= \textstyle{t + \log((1 + \lambda)/(\exp[2t] + \lambda)) } \eqsp ,
    \end{align}
    which concludes the proof.
  \end{proof}


  \begin{lemma}
  \label{lemma:integral_exp_uno}
  Let $\lambda \in \ooint{-1, +\infty}$ and $a: \ \ccint{0,T} \to \rset$ such
  that for any $t \in \ccint{0,T}$,
  \begin{equation}
    a(t) = 1 - 2/(1 + \exp[-2(T-t)]\lambda) \eqsp . 
  \end{equation}
  Then,  we have that for any $t \in \ccint{0,T}$,
  \begin{align}
      \textstyle{\int_0^t \exp[ 2\int_{T-s}^T a(u) \rmd u] \rmd s} &=     \textstyle{(1/2)(1+\lambda)^2 [(1 + \lambda \exp[-2t])^{-1} - 1/(1+\lambda)]/\lambda} \\
      &= \textstyle{(1/2)(1+\lambda)  (1 -  \exp[-2t])/(1 + \lambda \exp[-2t])  } \eqsp . 
    \end{align}
  \end{lemma}

  \begin{proof}
    Let $t \in \ccint{0,T}$. Using \Cref{lemma:integral_uno} we have that for any $s \in \ccint{0,T}$
    \begin{equation}
      \textstyle{ \exp[2\int_{T-s}^T a(u) \rmd u] = (1+\lambda)^2 \exp[2s] / (\lambda + \exp[2s])^2 = (1+\lambda)^2 \exp[-2s]/(1 + \lambda \exp[-2s])^2 \eqsp .} 
    \end{equation}
    Assume that $\lambda \neq 0$. Then, we have that
    \begin{align}
      \textstyle{\int_0^t \exp[2\int_{T-s}^T a(u) \rmd u] \rmd s } &= \textstyle{(1/2)(1+\lambda)^2/\lambda\int_0^t 2 \lambda \exp[-2t]/(1 + \lambda \exp[-2t])^2 \rmd s} \\
                                                                  &= \textstyle{(1/2)(1+\lambda)^2 [(1 + \lambda \exp[-2t])^{-1} - 1/(1+\lambda)]/\lambda  } \\
      &= \textstyle{(1/2)(1+\lambda)  (1 -  \exp[-2t])/(1 + \lambda \exp[-2t])  } \eqsp .
    \end{align}
    We conclude the proof upon remarking that his result still holds in the case where $\lambda =0$.
  \end{proof}

  \begin{lemma}
  \label{lemma:integral_exp_duo}
  Let $\lambda \in \ooint{-1, +\infty}$ and $a: \ \ccint{0,T} \to \rset$ such
  that for any $t \in \ccint{0,T}$,
  \begin{equation}
    a(t) = 1 - 2/(1 + \exp[-2(T-t)]\lambda) \eqsp . 
  \end{equation}
  Then, if $\lambda \neq 0$, we have that for any $t \in \ccint{0,T}$,
  \begin{align}
    &\textstyle{
      \int_0^t \exp[ 2\int_{T-s}^T a(u) \rmd u]/(1 + \lambda \exp[-2s]) \rmd s} \\
    & \qquad \qquad \qquad \textstyle{= (1/4)(1+\lambda)^2 [(1 + \lambda \exp[-2t])^{-2} - 1/(1+\lambda)^2]/\lambda 
      } \\
    & \qquad \qquad \qquad \textstyle{=   (1/4)(1 - \exp[-2t])(2 + \lambda (1 + \exp[-2t]))/(1 + \lambda \exp[-2t])^2 } \eqsp . 
    \end{align}
    If $\lambda =0$ we have
    \begin{equation}
      \label{eq:2}
      \textstyle{
        \int_0^t \exp[ 2\int_{T-s}^T a(u) \rmd u]/(1 + \lambda \exp[-2s]) \rmd s = (1/2)(1 - \exp[-2t]) \eqsp . 
        }
    \end{equation}

  \end{lemma}

  \begin{proof}
    Let $t \in \ccint{0,T}$. Using \Cref{lemma:integral_uno} we have that for any $s \in \ccint{0,T}$
    \begin{equation}
      \textstyle{ \exp[\int_{T-s}^T a(u) \rmd u]/(1 + \lambda \exp[-2s]) = (1+\lambda)^2 \exp[-2s] / (1 + \lambda \exp[-2s])^3 \eqsp .} 
    \end{equation}
    Assume that $\lambda \neq 0$. Then, we have that
    \begin{align}
      &\textstyle{\int_0^t \exp[\int_{T-s}^T a(u) \rmd u]/(1 + \lambda \exp[-2s]) \rmd s } = \textstyle{(1/2)(1+\lambda)^2/\lambda\int_0^t 2 \lambda \exp[-2t]/(1 + \lambda \exp[-2t])^3 \rmd s} \\
                                                                  & \qquad \qquad = \textstyle{(1/4)(1+\lambda)^2 [(1 + \lambda \exp[-2t])^{-2} - 1/(1+\lambda)^2]/\lambda  } \\
      & \qquad \qquad = \textstyle{  (1/4)(1 - \exp[-2t])(2 + \lambda (1 + \exp[-2t]))/(1 + \lambda \exp[-2t])^2  } \eqsp .
    \end{align}
    We conclude the proof upon remarking that his result still holds in the case where $\lambda =0$.
  \end{proof}

  \begin{lemma}
    \label{lemma:hola_pas_carre}
  Let $\lambda \in \ooint{-1, +\infty}$ and $a: \ \ccint{0,T} \to \rset$ such
  that for any $t \in \ccint{0,T}$,
  \begin{equation}
    a(t) = 1 - 2/(1 + \exp[-2(T-t)]\lambda) \eqsp . 
  \end{equation}
  Then, we have that for any $t \in \ccint{0,T}$
  \begin{align}
    &\textstyle{
      \int_0^t \exp[ 2\int_{T-s}^T a(u) \rmd u] a(T-s) \rmd s} \\
    & \qquad \qquad \textstyle{= -(1/2)(1- \exp[-2t])(1 - \lambda^2 \exp[-2t])/(1 + \lambda \exp[-2t])^2 \eqsp . 
      }
  \end{align}
  % If $\lambda = 0$, we have
  % \begin{equation}
  %   \textstyle{ \int_0^t \exp[ 2\int_{T-s}^T a(u) \rmd u] a(T-s) \rmd s = -(1/2)(1 - \exp[-2t]) \eqsp .}
  % \end{equation}  
\end{lemma}

\begin{proof}
  Let $t \in \ccint{0,T}$. We have that
  \begin{align}
    &\textstyle{\int_0^t \exp[ 2\int_{T-s}^T a(u) \rmd u] a(T-s) \rmd s} \\
    & \qquad \textstyle{= \int_0^t \exp[ 2\int_{T-s}^T a(u) \rmd u] \rmd s - 2 \int_0^t  \exp[ 2\int_{T-s}^T a(u) \rmd u] / (1 + \lambda \exp[-2s]) \rmd s } \eqsp . \label{eq:combi_i}
  \end{align}
  Using \Cref{lemma:integral_exp_uno}, we have that
  \begin{equation}
    \textstyle{
      \int_0^t \exp[ 2\int_{T-s}^T a(u) \rmd u] \rmd t = (1/2)(1+\lambda)(1 - \exp[-2t])/(1 + \lambda \exp[-2t])\eqsp . 
    }
    \label{eq:uno_unO}
  \end{equation}
  In addition, using \Cref{lemma:integral_exp_duo}, we have 
  \begin{align}
    &\textstyle{
      \int_0^t  \exp[ 2\int_{T-s}^T a(u) \rmd u] / (1 + \lambda \exp[-2s]) \rmd s} \\
    & \qquad \qquad \textstyle{= (1/4)(1 - \exp[-2t])(2 + \lambda(1 + \exp[-2t]))/(1 + \lambda \exp[-2t])^2 \eqsp . 
      } \label{eq:uno_duO}
  \end{align}
  Combining \eqref{eq:uno_unO} and \eqref{eq:uno_duO} in \eqref{eq:combi_i} we have that
  \begin{align}
    &\textstyle{\int_0^t \exp[ 2\int_{T-s}^T a(u) \rmd u] a(T-s) \rmd s} \\
    & \qquad = (1/2)(1 - \exp[-2t])[(1+\lambda)(1 + \lambda \exp[-2t])]/(1 + \lambda \exp[-2t])^2 \\
    & \qquad \qquad - (1/2)(1-\exp[-2t]) (2 + \lambda(1 + \exp[-2t]))/(1 + \lambda \exp[-2t])^2 \\
    & \qquad = -(1/2)(1- \exp[-2t])(1 - \lambda^2 \exp[-2t])/(1 + \lambda \exp[-2t])^2 \eqsp ,
  \end{align}
  which concludes the proof.
\end{proof}

  \begin{lemma}
    \label{lemma:integral_changement}
    Let $\lambda \in \ooint{-1, +\infty}$ we have that for any $t \in \ccint{0,T}$
    \begin{equation}
      \textstyle{\int_0^t (1 + \lambda \exp[-2s])^{-1}} \rmd s = (1/2) \log((\lambda + \exp[2t])/(\lambda +1)) \eqsp . 
    \end{equation}
    In addition, we have for any $t \in \ccint{0,T}$
        \begin{equation}
      \textstyle{\int_0^t (1 + \lambda \exp[-2s])^{-2}} \rmd s = (1/2) \log((\lambda + \exp[2t])/(\lambda +1)) + (\lambda/2)[(\exp[2t] + \lambda)^{-1} - (\lambda + 1)^{-1}] \eqsp .
    \end{equation}
    Finally, we have that for any $t \in \ccint{0,T}$
        \begin{align}
      \textstyle{\int_0^t (1 + \lambda \exp[-2s])^{-3}} \rmd s &=  \textstyle{(1/2) \log((\lambda + \exp[2t])/(\lambda +1)) + \lambda [(\exp[2t] + \lambda)^{-1} - (\lambda + 1)^{-1}] } \\
      & \qquad \qquad - (\lambda^2/4)[(\exp[2t] + \lambda)^{-2} - (\lambda + 1)^{-2}] \eqsp . 
    \end{align}    
  \end{lemma}

  \begin{proof}
    Let $k \in \{1, 2, 3\}$. Using the change of variable $u \mapsto \exp[2u]$
    we have that
    \begin{equation}
      \textstyle{\int_0^{t}(1 + \lambda \exp[-2s])^{-k} \rmd s = (1/2) \int_1^{\exp[2t]} u^{k-1} / (u + \lambda) \rmd u \eqsp . }
    \end{equation}
    Therefore, we have that
    \begin{equation}
      \textstyle{\int_0^{t}(1 + \lambda \exp[-2s])^{-1} \rmd s = (1/2) \int_1^{\exp[2t]}  (u + \lambda)^{-1} \rmd u = (1/2)\log((\lambda + \exp[2t])/(\lambda +1)) \eqsp . }
    \end{equation}
    In addition, using that for any $u \in \ccint{0,T}$, $u = (u + \lambda) - \lambda$ we have that
    \begin{align}
      &\textstyle{\int_0^{t}(1 + \lambda \exp[-2s])^{-2} \rmd s} = \textstyle{(1/2) \int_1^{\exp[2t]}  u(u + \lambda)^{-2} \rmd u } \\
      & \qquad = \textstyle{(1/2) \int_1^{\exp[2t]}  (u + \lambda)^{-1} \rmd u - (\lambda/2) \int_1^{\exp[2t]} (u + \lambda)^{-2} \rmd u } \\
      & \qquad = \textstyle{(1/2)\log((\lambda + \exp[2t])/(\lambda +1)) + (\lambda/2)[(\exp[2t] + \lambda)^{-1} - (\lambda + 1)^{-1}] } \eqsp . 
    \end{align}
    Finally, using that for any $u \in \ccint{0,T}$, $u \in \ccint{0,T}$, $u^2 = (u + \lambda)^2 - 2\lambda(u+\lambda) + \lambda^2$ we have that
    \begin{align}
      &\textstyle{\int_0^{t}(1 + \lambda \exp[-2s])^{-3} \rmd s} = \textstyle{(1/2) \int_1^{\exp[2t]}  u^2(u + \lambda)^{-2} \rmd u } \\
      & \qquad = \textstyle{(1/2) \int_1^{\exp[2t]}  (u + \lambda)^{-1} \rmd u - \lambda \int_1^{\exp[2t]} (u + \lambda)^{-2} \rmd u + (\lambda^2/2) \int_1^{\exp[2t]} (u + \lambda)^{-3} \rmd u} \\
      & \qquad = \textstyle{(1/2) \log((\lambda + \exp[2t])/(\lambda +1)) + \lambda [(\exp[2t] + \lambda)^{-1} - (\lambda + 1)^{-1}] } \\
      & \qquad \qquad - (\lambda^2/4)[(\exp[2t] + \lambda)^{-2} - (\lambda + 1)^{-2}] \eqsp ,
    \end{align}
    which concludes the proof.
  \end{proof}
  
\begin{lemma}
  \label{lemma:integral_duo}
  Let $\lambda \in \ooint{-1, +\infty}$ and $a: \ \ccint{0,T} \to \rset$ such
  that for any $t \in \ccint{0,T}$,
  \begin{equation}
    a(t) = 1 - 2/(1 + \exp[-2(T-t)]\lambda) \eqsp . 
  \end{equation}
  Then, we have that for any $t \in \ccint{0,T}$,  
  \begin{equation}
    \textstyle{\int_{T-t}^T a(s)^2 \rmd s} = t - 2\lambda(1 - \exp[-2t]) / [(1 + \lambda)(1 + \lambda \exp[-2t])] \eqsp . 
  \end{equation}
  \end{lemma}
  
  \begin{proof}
    Let $t \in \ccint{0,T}$. Similarly to the proof of
    \Cref{lemma:integral_uno}, we have that
    $\int_{T-t}^T a(s) \rmd s = \int_0^{t} a(T-s) \rmd s$.  Define $b$ such that
    for any $t \in \ccint{0,T}$, $b(t) = a (T-t)$. In particular, we have that
    for any $t \in \ccint{0,T}$
    \begin{equation}
      b(t) = 1 - 2 / (1 +\lambda \exp[-2t]) \eqsp . 
    \end{equation}
    We have that
    \begin{equation}
      \textstyle{\int_{T-t}^Ta(s)^2 \rmd s = \int_0^t b(s)^2 \rmd s  = \int_0^t (1 - 4/(1 + \lambda \exp[-2s]) + 4/(1 + \lambda \exp[-2s])^2) \rmd s  \eqsp . }
    \end{equation}
    Combining this result and \Cref{lemma:integral_changement}, we have
  \begin{align}
    \textstyle{\int_{T-t}^T a(s)^2 \rmd s} &= t + 2\lambda [(\lambda + \exp[2t])^{-1} - (\lambda + 1)^{-1} ] \\
    &= t - 2\lambda(1 - \exp[-2t]) / [(1 + \lambda)(1 + \lambda \exp[-2t])] \eqsp . 
  \end{align}    
    \end{proof}

    \begin{lemma}
      \label{lemma:hola_carre}
  Let $\lambda \in \ooint{-1, +\infty}$ and $a: \ \ccint{0,T} \to \rset$ such
  that for any $t \in \ccint{0,T}$,
  \begin{equation}
    a(t) = 1 - 2/(1 + \exp[-2(T-t)]\lambda) \eqsp . 
  \end{equation}
  Then, if $\lambda \neq 0$, we have that 
  \begin{align}
    &\textstyle{\int_0^T \exp[2\int_{T-t}^T a(s) \rmd s] (\int_{T-t}^T a(s)^2 \rmd s) \rmd t} = - (T/2)(1+\lambda)^2\exp[-2T]/(1+\lambda\exp[-2T]) \\
                                                                                             & \qquad \qquad + (1+\lambda)^2/(4\lambda)\log((1+\lambda)/(1 + \lambda \exp[-2T])) \\
    & \qquad \qquad -(\lambda/2)(1-\exp[-2T])^2/(1 + \lambda \exp[-2T])^2 \eqsp . \label{eq:big_eq_err}
  \end{align}
  If $\lambda = 0$, we have that
  \begin{equation}
    \textstyle{\int_0^T \exp[2\int_{T-t}^T a(s) \rmd s] (\int_{T-t}^T a(s)^2 \rmd s) \rmd t} = -(T/2) \exp[-2T] + (1/4)(1 - \exp[-2T]) \eqsp . \label{eq:small_eq}
  \end{equation}
\end{lemma}
Note that taking $\lambda \to 0$ in \eqref{eq:big_eq_err} we recover
\eqref{eq:small_eq}, using that for any $u > 0$,
$\lim_{\lambda \to 0} \log(1 + \lambda u)/\lambda = u$.

\begin{proof}
  We first start with the case $\lambda \neq 0$.
  Similarly to the proof of \Cref{lemma:integral_uno}, we have that
  $\int_{T-t}^T a(s) \rmd s = \int_0^{t} a(T-s) \rmd s$.  Define $b$ such that
  for any $t \in \ccint{0,T}$, $b(t) = a (T-t)$. We have that
  \begin{equation}
    \textstyle{
      \int_0^T \exp[2\int_{T-t}^T a(s) \rmd s] (\int_{T-t}^T a(s)^2 \rmd s) \rmd t = \int_0^T \exp[2\int_{T-t}^T a(s) \rmd s] (\int_{0}^t b(s)^2 \rmd s) \rmd t \eqsp .
      }
    \end{equation}
    Let $A: \ \ccint{0,T} \to \rset$ such that for any $t \in \ccint{0,T}$,
    \begin{equation}
      \textstyle{
        A(t) = \int_0^t \exp[2\int_{T-s}^T a(u) \rmd u] \rmd s \eqsp .
        }
    \end{equation}
    Note that $A(0) = 0$. Hence, by integration by parts, we have
    \begin{equation}
      \textstyle{ \int_0^T \exp[2\int_{T-t}^T a(s) \rmd s] (\int_{0}^t b(s)^2 \rmd s) \rmd t = A(T)\int_0^T b(t)^2 \rmd t - \int_0^T A(t) b(t)^2 \rmd t \eqsp . }
    \end{equation}
    In what follows, we compute $\int_0^T A(t) b(t)^2 \rmd t$. First, we recall that for any $t \in \ccint{0,T}$
    \begin{equation}\label{eq:b_square}
      b(t)^2 = (1 - 2/(1 + \lambda \exp[-2t]))^2 = 1 - 4/(1 + \lambda \exp[-2t]) + 4/(1 + \lambda \exp[-2t])^2 \eqsp . 
    \end{equation}
    In addition, using \Cref{lemma:integral_exp_uno}, we have that for any $t \in \ccint{0,T}$
    \begin{equation}
      \label{eq:A_exp}
      A(t) = (1/2)\{ (1+\lambda)^2/(\lambda(1 + \lambda \exp[-2t])) - (1+\lambda)/\lambda\} \eqsp . 
    \end{equation}
    Using \eqref{eq:b_square} and \eqref{eq:A_exp} we have that for any $t \in \ccint{0,T}$
    \begin{align}
      2A(t)b(t)^2 &= -(1+\lambda)/\lambda + [4(1+\lambda)/\lambda +(1+\lambda)^2/\lambda] u_1(t) \\
                 & \qquad - [4(1+\lambda)^2/\lambda + 4(1+\lambda)/\lambda] u_2(t) + [4(1+\lambda)^2/\lambda] u_3(t) \\
      &= -(1+\lambda)/\lambda + [(1+\lambda)(5 + \lambda)/\lambda] u_1(t) \\
      & \qquad - [4(1+\lambda)(2+\lambda)/\lambda] u_2(t) + [4(1+\lambda)^2/\lambda] u_3(t)\eqsp , \label{eq:dev_long}
    \end{align}
    where for any $k \in \{1,2, 3\}$ and $t \in \ccint{0,T}$ we have
    \begin{equation}
      u_k(t) = (1 + \lambda \exp[-2t])^{-k} \eqsp . 
    \end{equation}
    For any $k \in \{0, 1, 2\}$ denote $v_k : \ \ccint{0,T} \to \rset$ such that
    for any $t \in \ccint{0,T}$ and $k \in \{1, 2\}$
    \begin{equation}
      v_0(t) = \log((\lambda + \exp[2t])/(\lambda+1)) \eqsp , \qquad v_k(t) = (\exp[2t] + \lambda)^{-k} - (1 + \lambda)^{-k}\eqsp .
    \end{equation}
    Combining \eqref{eq:dev_long} and \Cref{lemma:integral_changement}, we get
    that for any $t \in \ccint{0,T}$
    \begin{align}
      &\textstyle{2 \int_0^t A(s) b(s)^2 \rmd s} = -[(1+\lambda)/\lambda]t \\
                                              & \qquad + (1/2) \{ [(1+\lambda)(5 + \lambda)/\lambda] - [4(1+\lambda)(2+\lambda)/\lambda] + [4(1+\lambda)^2/\lambda] \} v_0(t) \\
                                              & \qquad + \{ - (\lambda/2) [4(1+\lambda)(2+\lambda)/\lambda] + \lambda [4(1+\lambda)^2/\lambda]\} v_1(t) \\
                                              & \qquad  - (\lambda^2/4)  [4(1+\lambda)^2/\lambda] v_2(t) \\
      &= -[(1+\lambda)/\lambda]t + (1+\lambda)^2/(2\lambda) v_0(t) + 2(1+\lambda)\lambda v_1(t) - \lambda(1+\lambda)^2 v_2(t) \eqsp . 
    \end{align}
    In addition, we have that
    \begin{align}
      &-[(1+\lambda)/\lambda]t + (1+\lambda)^2/(2\lambda) v_0(t) = -[(1+\lambda)/\lambda]t + [(1+\lambda)^2/\lambda]t \\ 
                                                                & \qquad \qquad \qquad + (1+\lambda)^2/(2\lambda)\log((1 + \lambda \exp[-2t])/(1 + \lambda)) \\
      &\qquad \qquad =(1+\lambda)t +(1+\lambda)^2/(2\lambda)\log((1 + \lambda \exp[-2t])/(1 + \lambda)) \eqsp . 
    \end{align}
    Therefore, we get that
    \begin{align}
      \textstyle{2 \int_0^t A(s) b(s)^2 \rmd s} & = (1+\lambda)t +(1+\lambda)^2/(2\lambda)\log((1 + \lambda \exp[-2t])/(1 + \lambda))  \\
      & \qquad + 2(1+\lambda)\lambda v_1(t) + \lambda(1+\lambda)^2 v_2(t) \eqsp . \label{eq:et_un}
    \end{align}
    In addition, we have that
    \begin{equation}
      (1+\lambda)\lambda v_1(t) = -\lambda(1 - \exp[-2T])/(1 + \lambda \exp[-2T]) \eqsp . \label{eq:et_deux}
    \end{equation}
    We also have that
    \begin{align}
      \lambda(1+\lambda)^2 v_2(t) &= \lambda (2 \lambda + 1 - 2 \lambda \exp[2T] - \exp[4T])/(\exp[2T] + \lambda)^2 \\
                                  &= \lambda (1 - \exp[2T])(1 + 2\lambda  + \exp[2T]) /(\exp[2T] + \lambda)^2 \\
                                    &= -\lambda (1 - \exp[-2T])(1 + (1 + 2\lambda) \exp[-2T]) /(1 + \lambda \exp[-2T])^2 \eqsp . \label{eq:et_trois}
    \end{align}
    Finally, using \Cref{lemma:integral_exp_uno} and \Cref{lemma:integral_duo} we have
    \begin{align}
      \textstyle{A(T) \int_0^T b(t)^2 \rmd t} &= \textstyle{(1/2)(1+\lambda)  (1 -  \exp[-2T])/(1 + \lambda \exp[-2T])}\\
                                              & \qquad \times (T - 2 \lambda(1 - \exp[-2T]) / [(1 + \lambda)(1 + \lambda \exp[-2T])])   \\
                                              & = (T/2)(1+\lambda)  (1 -  \exp[-2T])/(1 + \lambda \exp[-2T]) \\
      & \qquad - \lambda (1 -  \exp[-2T])^2/(1 + \lambda \exp[-2T])^2 \eqsp .  \label{eq:et_quatre}
    \end{align}
    Combining \eqref{eq:et_un}, \eqref{eq:et_deux},
    \eqref{eq:et_trois} and \eqref{eq:et_quatre} we get
  \begin{align}
    &\textstyle{\int_0^T \exp[2\int_{T-t}^T a(s) \rmd s] (\int_{T-t}^T a(s)^2 \rmd s) \rmd t} = (T/2)(1+\lambda)(1-\exp[-2T])/(1+\lambda \exp[-2T]) \\
                                                                                             & \qquad \qquad - \lambda(1- \exp[-2T])^2/(1 + \lambda \exp[-2T])^2 \\
                                                                                             & \qquad \qquad - (1 + \lambda)(T/2) + (1+\lambda)^2/(4\lambda)\log((1+\lambda)/(1 + \lambda \exp[-2T])) \\
    & \qquad \qquad +\lambda(1 - \exp[-2T])/(1 + \lambda \exp[-2T]) \\
    & \qquad \qquad -(\lambda/2)(1-\exp[-2T])((1+2\lambda)\exp[-2T] +1)/(1 + \lambda \exp[-2T])^2 \eqsp . \label{eq:big_eq_err_inter}
  \end{align}
  In addition, we have that
  \begin{align}
    &-(\lambda/2)(1-\exp[-2T])^2/(1 + \lambda \exp[-2T])^2 & \\
    &\qquad \qquad \qquad = - \lambda(1- \exp[-2T])^2/(1 + \lambda \exp[-2T])^2 \\
                                                                            & \qquad \qquad \qquad \qquad +\lambda(1 - \exp[-2T])/(1 + \lambda \exp[-2T]) \\
    & \qquad \qquad \qquad \qquad -(\lambda/2)(1-\exp[-2T])((1+2\lambda)\exp[-2T] +1)/(1 + \lambda \exp[-2T])^2 \eqsp .
  \end{align}
  Combining this result and \eqref{eq:big_eq_err_inter}, we get
  \begin{align}
    &\textstyle{\int_0^T \exp[2\int_{T-t}^T a(s) \rmd s] (\int_{T-t}^T a(s)^2 \rmd s) \rmd t} = (T/2)(1+\lambda)(1-\exp[-2T])/(1+\lambda \exp[-2T]) \\
                                                                                             & \qquad \qquad - (1 + \lambda)(T/2) + (1+\lambda)^2/(4\lambda)\log((1+\lambda)/(1 + \lambda \exp[-2T])) \\
    & \qquad \qquad -(1/2)(1-\exp[-2T])^2/(1 + \lambda \exp[-2T])^2 \eqsp . \label{eq:big_eq_err_inter2}
  \end{align}
  Finally, we have
  \begin{align}
    &(T/2)(1+\lambda)(1-\exp[-2T])/(1+\lambda \exp[-2T]) - (T/2)(1+\lambda) \\
    &\qquad = - (T/2)(1+\lambda)^2\exp[-2T]/(1+\lambda\exp[-2T]) \eqsp ,
  \end{align}
  which concludes the proof in the case $\lambda \neq 0$ upon combining this
  result and \eqref{eq:big_eq_err_inter2}. In the case $\lambda = 0$, we have that
  for any $t \in \ccint{0,T}$, $a(t) = -1$ and therefore by integration by part
  we have
    \begin{equation}
      \textstyle{\int_0^T \exp[2\int_{T-t}^T a(s) \rmd s] (\int_{T-t}^T a(s)^2 \rmd s) \rmd t} = -(T/2) \exp[-2T] + (1/4)(1 - \exp[-2T]) \eqsp ,
    \end{equation}
    which concludes the proof.
    % Note that taking the limit $\lambda \to 0$ we have
    % \begin{equation}
    %   \textstyle{
    %     \lim_{\lambda \to 0} 2 \int_0^t A(s) b(s)^2 \rmd s = t + (1/2) (\exp[-2t] - 1) \eqsp .
    %     }
    %   \end{equation}
    %   In addition, using \Cref{lemma:integral_exp_uno} and \Cref{lemma:integral_duo} we have 
    %   \begin{equation}
    %     \textstyle{
    %       \lim_{\lambda \to 0} A(T) \int_0^T b(t)^2 \rmd t = (T/2) (1 - \exp[-2T])  \eqsp .
    %       }
    % \end{equation}
    % Hence, we have
    % \begin{equation}
    %   \textstyle{ \lim_{\lambda \to 0} \{ A(T)\int_0^T b(t)^2 \rmd t  -\int_0^T A(t) b(t)^2 \rmd t\} = -(T/2) \exp[-2T] +(1/4)(1 - \exp[-2T]) \eqsp . }
    % \end{equation}
  \end{proof}

    We are now ready to prove the following results.

    \begin{proposition}
      \label{prop:expression_fat_1}
      Let $\lambda \in \ooint{-1, +\infty}$ and $a: \ \ccint{0,T} \to \rset$ such
  that for any $t \in \ccint{0,T}$,
  \begin{equation}
    a(t) = 1 - 2/(1 + \exp[-2(T-t)]\lambda) \eqsp . 
  \end{equation}
  Then,  we have that for any $t \in \ccint{0,T}$,
  \begin{align}
    &\textstyle{\exp[2\int_0^T a(t) \rmd t]\{ \int_0^Ta(t)^2 \rmd t + a(T) - a(0) + 1\}} \\
    & \qquad \qquad \qquad = (T+1)\exp[-2T](\lambda+1)^2/(1 +\lambda \exp[-2T])^2 \eqsp . 
%    & \qquad \qquad \qquad \qquad -\lambda(1+\lambda)\exp[-2T](1-\exp[-2T])/(1 +\lambda \exp[-2T])^3 \eqsp . 
  \end{align}
  \end{proposition}

  \begin{proof}
    First, we have that
    \begin{align}
      a(T) - a(0) &= 1 - 2 /(1+\lambda) - 1 + 2/(1+\lambda \exp[-2T]) \\
                  &= 2\lambda(1 - \exp[-2T])/[(1+\lambda)(1+\lambda \exp[-2T])] \eqsp .
                    \label{eq:al_uno}
    \end{align}
    In addition, using \Cref{lemma:integral_duo} we have
  \begin{equation}
    \textstyle{\int_{0}^T a(s)^2 \rmd s} = T - 2 \lambda(1 - \exp[-2T]) / [(1 + \lambda)(1 + \lambda \exp[-2T])] \eqsp .
    \label{eq:al_duo}
  \end{equation}
  Finally, using \Cref{lemma:integral_uno} we have that
  \begin{equation}
          \textstyle{
      \exp[2\int_{0}^T a(s) \rmd s] = \exp[-2T] (\lambda+1)^2 /(1+ \lambda\exp[-2T])^2 \eqsp . 
      }\label{eq:al_tertio}
    \end{equation}
    We conclude the proof upon combining \eqref{eq:al_uno}, \eqref{eq:al_duo}
    and \eqref{eq:al_tertio}.
  \end{proof}

  Finally, we have the following proposition.

  \begin{proposition}
    \label{prop:expression_fat_2}
      Let $\lambda \in \ooint{-1, +\infty}$ and $a: \ \ccint{0,T} \to \rset$ such
  that for any $t \in \ccint{0,T}$,
  \begin{equation}
    a(t) = 1 - 2/(1 + \exp[-2(T-t)]\lambda) \eqsp . 
  \end{equation}
  Then, if $\lambda \neq 0$, we have that for any $t \in \ccint{0,T}$,
  \begin{align}
  &\textstyle{\int_0^T \exp[2\int_{T-t}^T a(s) \rmd s] [\int_{T-t}^T a(s)^2 \rmd s + a(T) - a(T-t)] \rmd t } \\
  & \qquad = - (T/2)(1+\lambda)^2\exp[-2T]/(1+\lambda\exp[-2T]) \\
   & \qquad \qquad + (1+\lambda)^2/(4\lambda)\log((1+\lambda)/(1 + \lambda \exp[-2T])) \\
    & \qquad \qquad + (\lambda/2) \exp[-2T]/(1 + \lambda \exp[-2T])^2 \eqsp .
  \end{align}
  If $\lambda = 0$, we have that
  \begin{align}
    &\textstyle{\int_0^T \exp[2\int_{T-t}^T a(s) \rmd s] [\int_{T-t}^T a(s)^2 \rmd s + a(T) - a(T-t)] \rmd t } \\
    & \qquad \qquad  \qquad \qquad = -(T/2)\exp[-2T] + (1/4)(1 - \exp[-2T]) \eqsp . 
  \end{align}

  \end{proposition}

  \begin{proof}
    We assume that $\lambda \neq 0$. The case where $\lambda = 0$ is left to the
    reader.  First, using \Cref{lemma:hola_carre}, we have that
    \begin{align}
          &\textstyle{\int_0^T \exp[2\int_{T-t}^T a(s) \rmd s] (\int_{T-t}^T a(s)^2 \rmd s) \rmd t} = - (T/2)(1+\lambda)^2\exp[-2T]/(1+\lambda\exp[-2T]) \\
                                                                                             & \qquad \qquad + (1+\lambda)^2/(4\lambda)\log((1+\lambda)/(1 + \lambda \exp[-2T])) \\
          & \qquad \qquad +(3\lambda/2)\exp[-2T](1-\exp[-2T])(1+\lambda)/(1 + \lambda \exp[-2T])^2 \eqsp .
              \label{eq:almost_uno}
    \end{align}
  Second, using \Cref{lemma:hola_pas_carre}, we have that
  \begin{align}
        &\textstyle{
      \int_0^T \exp[ 2\int_{T-t}^T a(u) \rmd u] a(T-t) \rmd t} \\
    & \qquad \qquad \textstyle{= -(1/2)(1- \exp[-2T])(1 - \lambda^2 \exp[-2T])/(1 + \lambda \exp[-2T])^2 \eqsp . \label{eq:almost_duo}
      }
  \end{align}
  Third, using \Cref{lemma:integral_exp_uno} and that $a(T) = 1 -2/(1+\lambda)$, we have that
  \begin{align}
    \textstyle{
    a(T) \int_0^T \exp[2\int_{T-t}^T a(s) \rmd s] \rmd t} &= a(T) \textstyle{(1/2)(1+\lambda)  (1 -  \exp[-2T])/(1 + \lambda \exp[-2T])  }\\
                                                          &=(1/2)(1+\lambda)  (1 -  \exp[-2T])/(1 + \lambda \exp[-2T]) \\
    & \qquad  -  (1 -  \exp[-2T])/(1 + \lambda \exp[-2T]) \eqsp .\label{eq:almost_tertio}
  \end{align}
  Combining \eqref{eq:almost_uno}, \eqref{eq:almost_duo} and \eqref{eq:almost_tertio} we get
\begin{align}
  &\textstyle{\int_0^T \exp[2\int_{T-t}^T a(s) \rmd s] [\int_{T-t}^T a(s)^2 \rmd s + a(T) - a(T-t)] \rmd t } \\
  & \qquad = - (T/2)(1+\lambda)^2\exp[-2T]/(1+\lambda\exp[-2T]) \\
   & \qquad \qquad + (1+\lambda)^2/(4\lambda)\log((1+\lambda)/(1 + \lambda \exp[-2T])) \\
   & \qquad \qquad +(3\lambda/2)\exp[-2T](1-\exp[-2T])(1+\lambda)/(1 + \lambda \exp[-2T])^2 \\
   & \qquad \qquad  +(1/2)(1- \exp[-2T])(1 - \lambda^2 \exp[-2T])/(1 + \lambda \exp[-2T])^2 \\ 
      & \qquad \qquad + (1/2)(1+\lambda)  (1 -  \exp[-2T])/(1 + \lambda \exp[-2T]) \\
     & \qquad \qquad  -  (1 -  \exp[-2T])/(1 + \lambda \exp[-2T])
\end{align}
In addition, we have
\begin{align}
  & (\lambda/2) (1 - \exp[-2T])/(1 + \lambda \exp[-2T])^2 \\
  &\qquad = (1/2)(1- \exp[-2T])(1 - \lambda^2 \exp[-2T])/(1 + \lambda \exp[-2T])^2 \\ 
      & \qquad \qquad + (1/2)(1+\lambda)  (1 -  \exp[-2T])/(1 + \lambda \exp[-2T]) \\
     & \qquad \qquad  -  (1 -  \exp[-2T])/(1 + \lambda \exp[-2T]) \eqsp ,
\end{align}
Finally, we have
\begin{align}
  &(\lambda/2)\exp[-2T]/(1 + \lambda \exp[-2T])^2 \\
    & \qquad \qquad = -(\lambda/2)(1-\exp[-2T])^2/(1 + \lambda \exp[-2T])^2 \\
    & \qquad \qquad \qquad + (\lambda/2) (1 - \exp[-2T])/(1 + \lambda \exp[-2T])^2 \eqsp .
\end{align}
which concludes the proof.
  % \begin{align}
  %   &\textstyle{\int_0^T \exp[2\int_{T-t}^T a(s) \rmd s] [\int_{T-t}^T a(s)^2 \rmd s + a(T) - a(T-t)] \rmd t } \\
  %     & \qquad = (T/2)(1+\lambda)(1-\exp[-2T])/(1+\lambda \exp[-2T]) \\
  %                                                                                            & \qquad \qquad - (\lambda/2)(1- \exp[-2T])^2/(1 + \lambda \exp[-2T])^2 \\
  %                                                                                            & \qquad \qquad - (1 + \lambda)(T/2) + (1+\lambda)^2/(4\lambda)\log((1+\lambda)/(1 + \lambda \exp[-2T])) \\
  %   & \qquad \qquad +(\lambda/2)(1 - \exp[-2T])/(1 + \lambda \exp[-2T]) \\
  %   & \qquad \qquad +(\lambda/3)(1-\exp[-2T])((1+2\lambda)\exp[-2T] +1)/(1 + \lambda \exp[-2T])^2 \\
  %   & \qquad \qquad +(1/2)(1- \exp[-2T])(1 - \lambda^2 \exp[-2T])/(1 + \lambda \exp[-2T])^2 \\
  %   & \qquad \qquad + (1/2)(1+\lambda)  (1 -  \exp[-2T])/(1 + \lambda \exp[-2T]) \\
  %   & \qquad \qquad - (1+\lambda)  (1 -  \exp[-2T])/(1 + \lambda \exp[-2T])^2 \\
  %     & \qquad = -(T/2)(1+\lambda)^2\exp[-2T]/(1+\lambda \exp[-2T]) \\
  %                                                                                            & \qquad \qquad - (\lambda/2)(1- \exp[-2T])^2/(1 + \lambda \exp[-2T])^2 \\
  %                                                                                            & \qquad \qquad  + (1+\lambda)^2/(4\lambda)\log((1+\lambda)/(1 + \lambda \exp[-2T])) \\
  %   & \qquad \qquad +(\lambda/2)(1 - \exp[-2T])/(1 + \lambda \exp[-2T]) \\
  %   & \qquad \qquad +(\lambda/3)(1-\exp[-2T])((1+2\lambda)\exp[-2T] +1)/(1 + \lambda \exp[-2T])^2 \\
  %   & \qquad \qquad +(1/2)(1- \exp[-2T])(1 - \lambda^2 \exp[-2T])/(1 + \lambda \exp[-2T])^2 \\
  %   & \qquad \qquad + (1/2)(1+\lambda)  (1 -  \exp[-2T])/(1 + \lambda \exp[-2T]) \\
  %   & \qquad \qquad - (1+\lambda)  (1 -  \exp[-2T])/(1 + \lambda \exp[-2T])^2 \eqsp .
  % \end{align}
  % In addition, we have that
  % \begin{equation}
  %   (1 - \lambda^2 \exp[-2T]) + (1+\lambda)(1 + \lambda \exp[-2T]) - 2(1+\lambda) = - \lambda(1 - \exp[-2T) \eqsp . 
  % \end{equation}
  % Therefore, we get that
  %   \begin{align}
  %     &\textstyle{\int_0^T \exp[2\int_{T-t}^T a(s) \rmd s] [\int_{T-t}^T a(s)^2 \rmd s + a(T) - a(T-t)] \rmd t } \\
  %     & \qquad = -(T/2)(1+\lambda)^2\exp[-2T]/(1+\lambda \exp[-2T]) \\
  %                                                                                            & \qquad \qquad - \lambda (1- \exp[-2T])^2/(1 + \lambda \exp[-2T])^2 \\
  %                                                                                            & \qquad \qquad  + (1+\lambda)^2/(4\lambda)\log((1+\lambda)/(1 + \lambda \exp[-2T])) \\
  %   & \qquad \qquad +(\lambda/2)(1 - \exp[-2T])/(1 + \lambda \exp[-2T]) \\
  %     & \qquad \qquad +(\lambda/3)(1-\exp[-2T])((1+2\lambda)\exp[-2T] +1)/(1 + \lambda \exp[-2T])^2 \eqsp ,
  %   \end{align}
  %   which concludes the proof
  \end{proof}

  Finally, we have the following result.

  \begin{proposition}
    \label{prop:explicit_lambda}
        Let $\lambda \in \ooint{-1, +\infty}$ and $a: \ \ccint{0,T} \to \rset$ such
  that for any $t \in \ccint{0,T}$,
  \begin{equation}
    a(t) = 1 - 2/(1 + \exp[-2(T-t)]\lambda) \eqsp . 
  \end{equation}
  Then, if $\lambda \neq 0$, we have that for any $t \in \ccint{0,T}$,
%   \begin{align}
%   &\textstyle{\int_0^T \exp[2\int_{T-t}^T a(s) \rmd s] [\int_{T-t}^T a(s)^2 \rmd s + a(T) - a(T-t)] \rmd t } \\
%   & \qquad = - (T/2)(1+\lambda)^2\exp[-2T]/(1+\lambda\exp[-2T]) \\
%    & \qquad \qquad + (1+\lambda)^2/(4\lambda)\log((1+\lambda)/(1 + \lambda \exp[-2T])) \\
%   & \qquad \qquad +(3\lambda/2)\exp[-2T](1-\exp[-2T])(1+\lambda)/(1 + \lambda \exp[-2T])^2 \\
%     & \qquad \qquad + (\lambda/2) (1 - \exp[-2T])/(1 + \lambda \exp[-2T])^2 \eqsp .
%   \end{align}
%   and
%     \begin{align}
%     &\textstyle{\exp[2\int_0^T a(t) \rmd t]\{ \int_0^Ta(t)^2 \rmd t + a(T) - a(0) + 1\}} \\
%     & \qquad \qquad \qquad = (T+1)\exp[-2T](\lambda+1)^2/(1 +\lambda \exp[-2T])^2 \eqsp . 
% %    & \qquad \qquad \qquad \qquad -\lambda(1+\lambda)\exp[-2T](1-\exp[-2T])/(1 +\lambda \exp[-2T])^3 \eqsp . 
%   \end{align}
  \begin{align}
    &\textstyle{-\exp[2\int_0^T a(t) \rmd t]\{ \int_0^Ta(t)^2 \rmd t + a(T) - a(0)\} + 1 - \exp[2\int_0^T a(t) \rmd t]} \\
    & \qquad \qquad - 2 \textstyle{\int_0^T \exp[2\int_{T-t}^T a(s) \rmd s] [\int_{T-t}^T a(s)^2 \rmd s + a(T) - a(T-t)] \rmd t } \\
    & \qquad \qquad \qquad = 1 - \exp[-2T] (1 - \lambda T \exp[-2T]) (\lambda+1)^2/(1 +\lambda \exp[-2T])^2 \\
   & \qquad \qquad \qquad \qquad - (1+\lambda)^2/(2\lambda)\log((1+\lambda)/(1 + \lambda \exp[-2T])) \\
  & \qquad \qquad \qquad \qquad -\lambda\exp[-2T]/(1 + \lambda \exp[-2T])^2 \eqsp . 
      \end{align}
% OLD OLD OLD  
%   \begin{align}
%     &\textstyle{-\exp[2\int_0^T a(t) \rmd t]\{ \int_0^Ta(t)^2 \rmd t + a(T) - a(0)\} + 1 - \exp[2\int_0^T a(t) \rmd t]} \\
%     & \qquad \qquad - 2 \textstyle{\int_0^T \exp[2\int_{T-t}^T a(s) \rmd s] [\int_{T-t}^T a(s)^2 \rmd s + a(T) - a(T-t)] \rmd t } \\
%     & \qquad \qquad \qquad = 1 - (1 + T)\exp[-2T](\lambda+1)^2/(1 +\lambda \exp[-2T])^2 \\
%      & \qquad \qquad \qquad\qquad + T(1+\lambda)^2\exp[-2T]/(1+\lambda \exp[-2T]) \\
%     & \qquad \qquad \qquad \qquad - (1+\lambda)^2/(2\lambda)\log((1+\lambda)/(1 + \lambda \exp[-2T])) \\
%     & \qquad \qquad \qquad \qquad +\lambda(1+\lambda)\exp[-2T](1-\exp[-2T])/(1 +\lambda \exp[-2T])^3 \\    
%                                                                                              & \qquad \qquad \qquad \qquad+ 2 \lambda (1- \exp[-2T])^2/(1 + \lambda \exp[-2T])^2 \\    
%     & \qquad \qquad \qquad \qquad -\lambda(1 - \exp[-2T])/(1 + \lambda \exp[-2T]) \\
%       & \qquad \qquad \qquad \qquad -(2\lambda/3)(1-\exp[-2T])((1+2\lambda)\exp[-2T] +1)/(1 + \lambda \exp[-2T])^2 \eqsp .
%   \end{align}
%   In particular, we have that
      In particular, we have that 
  \begin{align}
        & \textstyle{-\exp[2\int_0^T a(t) \rmd t]\{ \int_0^Ta(t)^2 \rmd t + a(T) - a(0)\} + 1 - \exp[2\int_0^T a(t) \rmd t]} \\
        & \qquad \qquad - 2 \textstyle{\int_0^T \exp[2\int_{T-t}^T a(s) \rmd s] [\int_{T-t}^T a(s)^2 \rmd s + a(T) - a(T-t)] \rmd t } \\
    & \qquad \qquad \qquad = 1 - (1/2) (1+\lambda)^2\log(1+\lambda)/\lambda + O(\exp[-2T]) \eqsp . 
  \end{align}  
  If $\lambda = 0$, we have that for any $t \in \ccint{0,T}$
  \begin{align}
        &\textstyle{-\exp[2\int_0^T a(t) \rmd t]\{ \int_0^Ta(t)^2 \rmd t + a(T) - a(0)\} + 1 - \exp[2\int_0^T a(t) \rmd t]} \\
        & \qquad \qquad - 2 \textstyle{\int_0^T \exp[2\int_{T-t}^T a(s) \rmd s] [\int_{T-t}^T a(s)^2 \rmd s + a(T) - a(T-t)] \rmd t } \\
    & \qquad \qquad \qquad = (1/2)(1 - \exp[-2T]) \eqsp . 
  \end{align}

  \end{proposition}

  \begin{proof}
    The proof is a direct consequence of \Cref{prop:expression_fat_1}, \Cref{prop:expression_fat_2} and the fact that
    \begin{align}
      &- \exp[-2T] (1 - \lambda T \exp[-2T]) (\lambda+1)^2/(1 +\lambda \exp[-2T])^2 \\      
      &\qquad =-(T+1)\exp[-2T] (1+\lambda^2)/(1+\lambda\exp[-2T])^2 \\
      &\qquad \qquad + T(1+\lambda)^2\exp[-2T]/(1+\lambda\exp[-2T]) \eqsp . 
    \end{align}
  \end{proof}

  \begin{proposition}
    \label{prop:mean_mean}
        Let $\lambda \in \ooint{-1, +\infty}$ and $a: \ \ccint{0,T} \to \rset$ such
  that for any $t \in \ccint{0,T}$,
  \begin{equation}
    a(t) = 1 - 2/(1 + \exp[-2(T-t)]\lambda) \eqsp . 
  \end{equation}
  Then, we have that
  \begin{align}
    &\textstyle{ \exp[2\int_0^T a(t) \rmd t] - 1 - \int_0^T\exp[2 \int_{T-t}^T a(s) \rmd s] \{ \int_{T-t}^T a(s)^2 \rmd s + a(T) - 3 a(T-t) \} \rmd t } \\
    & \qquad = \exp[-2T] (1+\lambda)^2 /(1+ \lambda\exp[-2T])^2 - 1 \\
    & \qquad \qquad -(1/2)(1+\lambda)  (1 -  \exp[-2T])/(1 + \lambda \exp[-2T])(1 -2/(1+\lambda)) \\
    & \qquad \qquad -(3/2)(1- \exp[-2T])(1 - \lambda^2 \exp[-2T])/(1 + \lambda \exp[-2T])^2 \\
& \qquad \qquad + (T/2)(1+\lambda)^2\exp[-2T]/(1+\lambda\exp[-2T]) \\
    & \qquad \qquad - (1+\lambda)^2/(4\lambda)\log((1+\lambda)/(1 + \lambda \exp[-2T])) \\
    & \qquad \qquad +(\lambda/2)(1-\exp[-2T])^2/(1 + \lambda \exp[-2T])^2 \eqsp .    
  \end{align}
  % OOOLLD OOOOLD
  % \begin{align}
  %   &\textstyle{ \exp[2\int_0^T a(t) \rmd t] - 1 - \int_0^T\exp[2 \int_{T-t}^T a(s) \rmd s] \{ \int_{T-t}^T a(s)^2 \rmd s + a(T) - 3 a(T-t) \} \rmd t } \\
  %   & \qquad = \exp[-2T] (1+\lambda)^2 /(1+ \lambda\exp[-2T])^2 - 1 \\
  %   & \qquad \qquad  - (T/2)(1+\lambda)(1-\exp[-2T])/(1+\lambda \exp[-2T]) \\    
  %   & \qquad \qquad + (\lambda/2)(1- \exp[-2T])^2/(1 + \lambda \exp[-2T])^2 \\
  %                                                                                            & \qquad \qquad + (1 + \lambda)(T/2) + (1+\lambda)^2/(4\lambda)\log((1+\lambda)/(1 + \lambda \exp[-2T])) \\
  %   & \qquad \qquad -(\lambda/2)(1 - \exp[-2T])/(1 + \lambda \exp[-2T]) \\
  %   & \qquad \qquad -(\lambda/3)(1-\exp[-2T])((1+2\lambda)\exp[-2T] +1)/(1 + \lambda \exp[-2T])^2 \\
  %   & \qquad \qquad -(1/2)(1+\lambda)  (1 -  \exp[-2t])/(1 + \lambda \exp[-2t])a(T) \\
  %   & \qquad \qquad -(3/2)(1- \exp[-2T])(1 - \lambda^2 \exp[-2T])/(1 + \lambda \exp[-2T])^2 \eqsp . 
  % \end{align}
  In particular, we have
  \begin{align}
    &\textstyle{ \exp[2\int_0^T a(t) \rmd t] - 1 - \int_0^T\exp[2 \int_{T-t}^T a(s) \rmd s] \{ \int_{T-t}^T a(s)^2 \rmd s + a(T) - 3 a(T-t) \} \rmd t } \\
    & \qquad \qquad =  -2 - (1+\lambda)^2/(4\lambda)\log(1+\lambda) + O(\exp[-2T])  \eqsp . 
  \end{align}
\end{proposition}

\begin{proof}
  Using \Cref{lemma:integral_uno}, we have that
  \begin{equation}
    \textstyle{\exp[2\int_0^T a(t) \rmd t] = \exp[-2T] (1+\lambda)^2 /(1+ \lambda\exp[-2T])^2 \eqsp . } \label{eq:uno_exp}
  \end{equation}
  Using \Cref{lemma:integral_exp_uno}, we have
  \begin{equation}
    \textstyle{ \int_0^T \exp[2  \int_{T-t}^T a(s) \rmd s] \rmd t =  (1/2)(1+\lambda)  (1 -  \exp[-2t])/(1 + \lambda \exp[-2t]) }\eqsp .\label{eq:duo_exp}
  \end{equation}
  Using \Cref{lemma:hola_pas_carre}, we have 
  \begin{align}
    &\textstyle{
      \int_0^T \exp[ 2\int_{T-t}^T a(s) \rmd s] a(T-t) \rmd t} \\
    & \qquad \qquad \textstyle{= -(1/2)(1- \exp[-2T])(1 - \lambda^2 \exp[-2T])/(1 + \lambda \exp[-2T])^2 \eqsp . 
      } \label{eq:tertio_exp}
  \end{align}
  Finally, using \Cref{lemma:hola_carre}, we have
  \begin{align}
    &\textstyle{\int_0^T \exp[2\int_{T-t}^T a(s) \rmd s] (\int_{T-t}^T a(s)^2 \rmd s) \rmd t} = - (T/2)(1+\lambda)^2\exp[-2T]/(1+\lambda\exp[-2T]) \\
    & \qquad \qquad + (1+\lambda)^2/(4\lambda)\log((1+\lambda)/(1 + \lambda \exp[-2T])) \\
    & \qquad \qquad -(\lambda/2)(1-\exp[-2T])^2/(1 + \lambda \exp[-2T])^2 \eqsp . \label{eq:quatro_exp}
  \end{align}
  We conclude upon combining \eqref{eq:uno_exp}, \eqref{eq:duo_exp},
  \eqref{eq:tertio_exp}, \eqref{eq:quatro_exp} and that
  $a(T) = 1- 2/(1+\lambda)$.
\end{proof}


\subsection{General setting}
\label{sec:main-content}

In this section, we prove \Cref{thm:control_diffusion}. In order to compare our
results with \cite[Theorem 1]{debortoli2021diffusion}, we redefine a few
processes.  Let $p \in \Pens(\rset^d)$ be the target distribution. Consider the
Ornstein-Ulhenbeck forward dynamics $(x_t)_{t \in \ccint{0,T}}$ such that
$\rmd x_t = -x_t \rmd t + \sqrt{2} \rmd w_t$ and $x_0$ has distribution
$p_0$. We consider the backward chain $(X_k)_{k \in \{0, \dots, N\}}$ such that
for any $k \in \{0, \dots, N-1\}$,
\begin{equation}
  \label{eq:reverse_discrete}
  X_k = X_{k+1} + \gamma_{k+1} \{ X_{k+1} + 2 \nabla \log p_{t_{k+1}}(X_{k+1}) \} + \sqrt{2 \gamma_{k+1}} Z_{k+1}  ,
\end{equation}
with $\{Z_k\}_{k \in \nset}$ a family of i.i.d. Gaussian random variables with
zero mean and identity covariance matrix, $t_{k} = \sum_{\ell=1}^k \gamma_{\ell}$,
$\sum_{\ell=1}^{N} \gamma_{\ell} = T$ and $X_N$ has distribution
$p_0 = \mathrm{N}(0, \Id)$ (independent from $\{Z_k\}_{k \in \nset}$). Notice
that here we do not consider a score approximation in the recursion in order to
clarify our approximation results. We recall the following result from
\cite[Theorem 1]{debortoli2021diffusion}.

\begin{theorem}
  Assume that $p_0$ admits a bounded density (w.r.t. the Lebesgue measure)
  $p_0 \in \rmc^3(\rset^d, \ooint{0, +\infty})$  and that there
  exist $d_1, A_1, A_2, A_3 \geq 0$, $\beta_1, \beta_2, \beta_3 \in \nset$ and
  $\mtt_1 > 0$ such that for any $x \in \rset^d$ and $i \in \{1, 2, 3\}$
            \begin{equation}
              \textstyle{
              \normLigne{\nabla^i \log p_0(x)} \leq A_i(1 + \normLigne{x}^{\beta_i}) , \quad \langle \nabla \log p_0(x), x \rangle \leq -\mtt_1 \norm{x}^2 + d_1 \norm{x} ,}
          \end{equation}
          with $\beta_1 = 1$.  Then  there exist
          $B, C, D \geq0$ such that for any $N \in \nset$
          and $\{\gamma_k\}_{k=1}^N$ with $\gamma_k > 0$ for any
          $k \in \{1, \dots, N\}$ we have 
          \begin{equation}
            \label{eq:ineq_thm_1}
 \tvnormLigne{\mathcal{L}(X_0)-p_0} \leq   C \exp[D T] \sqrt{\gamma^\star} + B \exp[-T]  .          
          \end{equation}
          where $\gamma^\star = \sup_{k \in \{1, \dots, N\}} \gamma_k$ and
          $\mathcal{L}(X_0)$ is the distribution of $X_0$ given in \eqref{eq:reverse_discrete}.
        \end{theorem}

In the rest of this note we improve the theorem in the following way:
% \begin{enumerate}[wide, labelwidth=!, labelindent=0pt, label=(\alph*)]
\begin{enumerate}[label=(\alph*)]
\item We remove the exponential dependency w.r.t. the time in the first term of the RHS of \eqref{eq:ineq_thm_1}.
\item We provide explicit bounds $B, C, D \geq 0$ depending on the parameters of $p_0$.
\end{enumerate}

\begin{lemma}
  \label{lemma:control_growth}
  Assume
  \begin{equation}\label{params:thm2_app}
    \textstyle{
  \sup_{x,t} \normLigne{\nabla^2 \log p_t(x)} \leq
  \Ktt}~~\mbox{and}~~
  \textstyle{\normLigne{\partial_t \nabla \log p_t(x)} \leq \Mtt\, \rme^{-\alpha t}\,
  \normLigne{x}}.
  \end{equation}
  Then there exists $D \geq 0$ such that for any
  $x \in \rset^d$ and $t \in \ccint{0,T}$,
  $\normLigne{\nabla \log p_t(x)} \leq D(1 + \normLigne{x})$ with
  $D = \normLigne{\nabla \log p_0(0)} + \Ktt + C T$.
\end{lemma}

\begin{proof}
  Let $x \in \rset^d$ and $t \in \ccint{0,T}$. Since
  $(t,x) \mapsto \log p_t(x) \in \rmc^{2}(\ccint{0,T}\times \rset^d,
  \ooint{0,+\infty})$, we have that
  \begin{align}
    \textstyle{
      \nabla \log p_t(x)} &= \textstyle{\nabla \log p_0(x) + \int_0^t \partial_s \nabla \log p_s(x) \rmd s }\\
    &=  \textstyle{\nabla \log p_0(0) + \int_0^1 \nabla^2 \log p_{0}(ux)(x) \rmd u +  \int_0^t \partial_s \nabla \log p_s(x) \rmd s .
      }
  \end{align}
  Therefore,  we have that
  \begin{align}
    \normLigne{\nabla \log p_t(x)} &\leq \textstyle{\normLigne{\nabla \log p_0(0)} + \Ktt \normLigne{x} +  \int_0^t \normLigne{\partial_s \nabla \log p_s(x)} \rmd s } \\
                                   &\leq \textstyle{\normLigne{\nabla \log p_0(0)} + \Ktt \normLigne{x} +  \Mtt \sum_{k=0}^{N-1}(t_k - t_{k-1}) \exp[-\alpha t_k] \normLigne{x} } \\
                                   &\leq \textstyle{\normLigne{\nabla \log p_0(0)} + \Ktt \normLigne{x} +  \Mtt T \normLigne{x} } ,
  \end{align}
which concludes the proof.
\end{proof}

Note that in the previous proposition we can derive a tighter bound for $D$
which does not depend on the limiting time $T > 0$. However, we do not use the
bound $D>0$ in our quantitative result and therefore our simple bound suffices.

We also have the following useful lemma.

\begin{lemma}
  \label{item:bound_p_infty}
  Let $T \geq \log(2 \expeLigne{\normLigne{X_0}^2}) + \log(2)/2$ and assume that there exists $\eta > 0$ such that $aaa$ Then, we have
  \begin{equation}    
    \textstyle{
      \int_{\rset^d} p_\infty(x_T)^2 / p_T(x_T) \rmd x_T \leq \exp[4] + E_T \eqsp ,
      }
    \end{equation}
    with $E_T  \sim C\exp[-T]$ when $T \to +\infty$ and $C \geq 0$.
\end{lemma}

If $p_\infty$ satisfies the following $\Phi$-entropy inequality for any
$f: \ \rset_d \to \ooint{0,\infty}$ measurable
\begin{equation}
  \label{eq:functional_inequality}
  \textstyle{
    \int_{\rset^d} \normLigne{\nabla f(x)}^2/f(x)^3 p_\infty(x) \rmd x \leq C [\int_{\rset^d} (1/f(x)) p_\infty(x) \rmd x - 1/(\int_{\rset^d} f(x) p_\infty(x) \rmd x)] \eqsp ,
    }
  \end{equation}
  with $C \geq 0$. Then, we have as in  \cite[Proposition 7.6.1]{bakry:gentil:ledoux:2014}
  \begin{equation}
    \textstyle{
      \chi^2(p_\infty||p_t) = \int_{\rset^d} p_\infty^2(x)/p_T(x) \rmd x - 1 \leq \rme^{-Ct} \eqsp ,
      }
  \end{equation}
  which immediately concludes the proof of \Cref{item:bound_p_infty}. However,
  to the best of our knowledge, establishing \eqref{eq:functional_inequality}
  remains an open problem. Note that controlling $\chi^2(p_t||p_\infty)$ is much
  easier as the exponential decay of this divergence is linked with the
  Poincar\'e inequality which is satisfied in our Gaussian setting. In what
  follows, we consider another approach which relies on the structure of the
  Ornstein-Ulhenbeck transition kernel and provide non-tight upper bounds.
  
\begin{proof}
  Let $T \geq 0$, $\vareps >0$ and $x_T \in \rset^d$
  \begin{equation}
    \normLigne{x_T - \rme^{-T} x_0}^2 \leq (1 + \vareps) \normLigne{x_T}^2 + (1+1/\vareps) \rme^{-2T} \normLigne{x_0}^2 \eqsp . 
  \end{equation}
  Let $\vareps >0$ and $x_T \in \rset^d$, we have
  \begin{align}
    &\textstyle{
      p_T(x_T)^{-2} \leq \exp[(1+\vareps)/\sigma_T^2\normLigne{x_T}^2 ] }\\
    & \qquad \qquad \qquad \qquad \times \textstyle{(\int_{\rset^d} p(x_0) \exp[-\rme^{-2T}(1+1/\vareps)/(2\sigma_T^2)\normLigne{x_0}^2]\rmd x_0)^{-2} (2\uppi\sigma_T^2)^d  \eqsp .
      }
  \end{align}
  For any $x_T \in \rset^d$, we have
  \begin{align}
    p_\infty(x_T)^2/p_T(x_T) &\leq \exp[\{-1+(1+\vareps)/(2\sigma_T^2)\}\normLigne{x_T}^2 ](2\uppi/\sigma_T^2)^{-d/2} \\ & \qquad \qquad \textstyle{ (\int_{\rset^d} p(x_0) \exp[-\rme^{-2T}(1+1/\vareps)/(2\sigma_T^2)\normLigne{x_0}^2] \rmd x_0)^{-1}  \eqsp . }
  \end{align}
  In what follows, we set $\vareps = \rme^{-T}$. We have that
  \begin{equation}
    -1 + (1+\vareps)/(2\sigma_T^2) = (2\sigma_T^2)^{-1}(-2\sigma_T^2 + 1 + \vareps) = -(1 - 2\rme^{-T} + \vareps)/(2\sigma_T^2) = -(1 - \rme^{-T})/(2\sigma_T^2) \eqsp . 
  \end{equation}
  Therefore, we get that
  \begin{equation}
    \label{eq:UNO_de}
    \textstyle{
      \int_{\rset^d} \exp[\{-1+(1+\vareps)/(2\sigma_T^2)\}\normLigne{x_T}^2 ](2\uppi/\sigma_T^2)^{-d/2} \rmd x_T = (1 - \rme^{-T})^{-d/2} \eqsp .
      }
    \end{equation}
    In addition, we have that for any $R \geq 0$ using that $\sigma_T^2 \geq 1/2$ since $T \geq \log(2)/2$
    \begin{align}
      &\textstyle{
        \int_{\rset^d} p(x_0) \exp[-\rme^{-2T}(1+1/\vareps)/\sigma_T^2\normLigne{x_0}^2] \rmd x_0 } \\
      & \qquad \qquad  \geq \textstyle{\probaLigne{X_0 \in \cball{0}{R}} \exp[-\rme^{-2T}(1+1/\vareps)/\sigma_T^2R^2]} \\ 
        & \qquad \qquad  \geq \textstyle{\probaLigne{X_0 \in \cball{0}{R}} \exp[-4\rme^{-T}R^2]
        } \label{eq:inter_DUO}
    \end{align}
    Now let $R^2 = \rme^T$. We obtain
    \begin{equation}
      \textstyle{
        \int_{\rset^d} p(x_0) \exp[-\rme^{-2T}(1+1/\vareps)/\sigma_T^2\normLigne{x_0}^2] \rmd x_0 } \geq \probaLigne{X_0 \in \cball{0}{\rme^{T/2}}} \exp[-4] \eqsp . 
    \end{equation}
    In addition, using Markov inequality, we have
    \begin{equation}
      \probaLigne{X_0 \in \cball{0}{\rme^{T/2}}} = 1 - \probaLigne{\normLigne{X_0}^2 \geq \rme^T} \geq 1 - \expeLigne{\normLigne{X_0}^2} \rme^{-T} \geq 0 \eqsp . 
    \end{equation}
    Therefore, combining this result and \eqref{eq:inter_DUO}, we have
    \begin{equation}
      \label{eq:DUO_de}
      \textstyle{
        \int_{\rset^d} p(x_0) \exp[-\rme^{-2T}(1+1/\vareps)/\sigma_T^2\normLigne{x_0}^2] \rmd x_0 } \geq \exp[-4] (1 - \expeLigne{\normLigne{X_0}^2} \rme^{-T}) > 0 \eqsp . 
    \end{equation}
    We conclude upon combining \eqref{eq:UNO_de} and \eqref{eq:DUO_de}.
\end{proof}


We are now ready to state the following lemma.

\begin{lemma}
  \label{lemma:strong_solution_backward}
  There exists a unique strong solution to the SDE
  $\rmd y_t = \{ y_t + 2 \nabla \log p_{T-t}(y_t)\} \rmd t + \sqrt{2}
  \rmd w_t$ with initial condition $\mathcal{L}(y_0) = p_\infty$. In
  addition, we have that
  $\expeLigne{\sup_{t \in \ccint{0,T}}\normLigne{y_t}^\alpha}<+\infty$ for any 
  $\alpha >0$.
\end{lemma}

\begin{proof}
  Let $b: \ \ccint{0,T} \times \rset^d$ given for any $t \in \ccint{0,T}$ and
  $x \in \rset^d$ by $b(t,x) = x + 2 \nabla \log p_t(x)$. We have that
  $b \in \rmc^1(\ccint{0,T} \times \rset^d, \rset^d)$ and in particular is
  locally Lipschitz. In addition, using \Cref{lemma:control_growth} we have that
  for any $t \in \ccint{0,T}$ and $x \in \rset^d$,
  $\normLigne{b(t,x)} \leq (1 + D)\normLigne{x}$. Hence using \cite[Theorem 2.3,
  Theorem 3.1]{ikeda1989sto} and \cite[Theorem 2.1]{meyn1993criteria_iii} (with
  $V(x) = (1/2)\normLigne{x}^2$) there exists a unique strong solution to the
  SDE
  $\rmd y_t = \{ y_t + 2 \nabla \log p_{T-t}(y_t)\} \rmd t + \sqrt{2}
  \rmd w_t$ with initial condition $\mathcal{L}(y_0) = p_\infty$. Let
  $\alpha > 1$, then we have for any $t \in \ccint{0,T}$
  \begin{align}
    \textstyle{\sup_{s \in \ccint{0,t}} \normLigne{y_t}^\alpha} &\leq \textstyle{3^{\alpha-1} [\normLigne{y_0}^\alpha + t^{\alpha-1}(1+D)^\alpha\int_0^t \sup_{u \in \ccint{0,s}} \normLigne{y_u}^\alpha \rmd u + 2^{\alpha/2} \sup_{s \in \ccint{0,t}} \normLigne{w_u}^\alpha]} . 
  \end{align}
  Using that $\expeLigne{\sup_{s \in \ccint{0,T}} \normLigne{w_u}^\alpha}$
  and Gr\"onwall's lemma, we get that
  $\expeLigne{\textstyle{\sup_{t \in \ccint{0,T}}
      \normLigne{y_t}^\alpha}}<+\infty$ for any $\alpha >1$. The result is
  extended to any $\alpha >0$ since for any $\alpha \in \ocint{0,1}$ we have
  that
  \begin{equation}
    \textstyle{\expeLigne{\textstyle{\sup_{t \in \ccint{0,T}}
      \normLigne{y_t}^\alpha}} \leq \expeLigne{\textstyle{\sup_{t \in \ccint{0,T}}
      \normLigne{y_t}}}^\alpha <+\infty} .
  \end{equation}

\end{proof}

We are now ready to prove \Cref{thm:control_diffusion}.

\begin{proof}
  The beginning of the proof is similar to the one of \cite[Theorem
  1]{debortoli2021diffusion}.  For any $k \in \{1, \dots, N\}$, denote $\Rker_k$
  the Markov kernel such that for any $x \in \rset^d$, $\msa \in \mcb{\rset^d}$
  and $k \in \{0, \dots, N-1\}$ we have
    \begin{equation}
      \textstyle{\Rker_{k+1}(x, \msa) = (4 \uppi \gamma_{k+1})^{-d/2} \int_{\msa} \exp[-\norm{\tilde{x} - \Tnplusun(x)}^2/(4 \gamma_{k+1})] \rmd \tilde{x}  , }
    \end{equation}
    where for any $x \in \rset^d$,
    $\Tnplusun(x) = x + \gamma_{k+1} \defEnsLigne{x + 2 \nabla \log
      p_{t_{k+1}}(x)}$. Define for any $k_0, k_1 \in \{1, \dots, N\}$ with
    $k_1 \geq k_0$ $\Qker_{k_0, k_1} = \prod_{\ell = k_0}^{k_1}
    \Rker_{k_1 + k_0 - \ell}$. Finally, for ease of notation, we also define for any
    $k \in \{1, \dots, N\}$, $\Qker_k = \Qker_{k+1,N}$. Note that for any
    $k \in \{1, \dots, N\}$, $X_k$ has distribution $p_{\infty} \Qker_k$,
    where $p_{\infty}\in \Pens(\rset^d)$ with density w.r.t. the Lebesgue
    measure $\pref$. Let $\Pbb \in \Pens(\contspace)$ be the probability measure
    associated with the diffusion
    \begin{equation}
      \rmd x_t = - x_t \rmd t + \sqrt{2} \rmd w_t  , \quad x_0 \sim p_0  , 
    \end{equation}
    First, we have for any $\msa \in \mcb{\rset^d}$
    \begin{equation}
      p_0 \Pbb_{T|0} (\Pbb^R)_{T|0} (\msa) = \Pbb_T (\Pbb^R)_{T|0} (\msa) = (\Pbb^R)_0 (\Pbb^R)_{T|0} (\msa)  = (\Pbb^R)_T(\msa) = p_0(\msa)  . 
    \end{equation}
    Hence $p_0 = p_0 \Pbb_{T|0} (\Pbb^R)_{T|0}$.  Using this result % and
    % \Cref{lemma:data_processing_tv},
    we have
    \begin{align}
      \tvnormLigne{p_0 - p_{\infty} \Qker_0} &= \tvnormLigne{p_0 \Pbb_{T|0} (\Pbb^R)_{T|0} - p_{\infty} \Qker_0} \\
                                         &\leq \tvnormLigne{p_0 \Pbb_{T|0} (\Pbb^R)_{T|0} - p_{\infty} (\Pbb^R)_{T|0}} +\tvnormLigne{p_{\infty} (\Pbb^R)_{T|0} - p_{\infty} \Qker_0} \\
          &\leq \tvnormLigne{p_0 \Pbb_{T|0}- p_{\infty}} +\tvnormLigne{p_{\infty} (\Pbb^R)_{T|0} - p_{\infty} \Qker_0}  .                                  
    \end{align}
    Note that $\mathcal{L}(X_0)= p_{\infty} \Qker_0$ and therefore
    \begin{equation}
      \tvnormLigne{\mathcal{L}(X_0) - p_0} \leq \tvnormLigne{p_0 \Pbb_{T|0}- p_{\infty}} +\tvnormLigne{p_{\infty} (\Pbb^R)_{T|0} - p_{\infty} \Qker_0}  .
    \end{equation}
    We now bound each one of these terms.
    \begin{enumerate}[wide, labelwidth=!, labelindent=0pt, label=(\alph*)]
    \item First, we bound $\tvnormLigne{p_0 \Pbb_{T|0}- p_{\infty}}$. Using the 
      Pinsker inequality \cite[Equation 5.2.2]{bakry:gentil:ledoux:2014} we have
      that
      \begin{equation}
        \label{eq:pinsker}
        \tvnormLigne{p_0 \Pbb_{T|0}- p_{\infty}} \leq \sqrt{2} \KLLignesqrt{p_0 \Pbb_{T|0}}{p_{\infty}}  . 
      \end{equation}
      In addition, $p_\infty$ satisfies the log-Sobolev inequality with
      constant $C=1$, \cite{gross1975logarithmic}. Namely, for any
      $f \in \rmc^1(\rset^d, \ooint{0,+\infty})$ such that
      $f \in \rmL^1(p_\infty)$ and
      $\int_{\rset^d} \normLigne{\nabla \log f(x)}^2 f(x) \rmd p_\infty(x) <
      +\infty$ we have
      \begin{align}
        &\textstyle{
          \int_{\rset^d} f(x) \log f(x) \rmd p_{\infty}(x) - (\int_{\rset^d} f(x) \rmd p_{\infty}(x)) (\log \int_{\rset^d} f(x) \rmd p_{\infty}(x)) } \\
        & \qquad \qquad \qquad \qquad  \textstyle{\leq (C/2) \int_{\rset^d}  \normLigne{\nabla \log f(x)}^2 f(x) \rmd p_\infty(x)  ,          
          }
      \end{align}
      with $C=1$. Therefore, using \cite[Theorem
      5.2.1]{bakry:gentil:ledoux:2014} we have that for any
      $f \in \rmL^1(p_\infty)$ with
      $\int_{\rset^d} \abs{f(x)} \abs{\log f(x)} \rmd p_\infty(x) < +\infty$      
      \begin{equation}
        \label{eq:entropy_decay}
        \Ent{p_\infty}{\Pbb_{T|0}[f]} \leq \exp[-2T]  \Ent{p_\infty}{f}  ,
      \end{equation}
      where for any $g \in \rmL^1(p_\infty)$ with 
      $\int_{\rset^d} \abs{g(x)} \abs{\log g(x)} \rmd p_\infty(x) < +\infty$ we define
      \begin{equation}
        \textstyle{
          \Ent{p_\infty}{g} = \int_{\rset^d} g(x) \log g(x) \rmd p_{\infty}(x) - (\int_{\rset^d} g(x) \rmd p_{\infty}(x)) (\log \int_{\rset^d} g(x) \rmd p_{\infty}(x)) .
          }
        \end{equation}
        Note that
        $(\rmd p_T / \rmd p_\infty) = \Pbb_{T|0} [\rmd p_0 / \rmd
        p_\infty]$ and that for any $\mu \in \Pens(\rset^d)$ with
        $\KL{\mu}{p_\infty} < +\infty$ we have
        $\Ent{p_\infty}{\rmd \mu / \rmd p_\infty} =
        \KLLigne{\mu}{p_\infty}$. Using these results, \eqref{eq:pinsker} and 
        \eqref{eq:entropy_decay} we get that
        \begin{equation}
          \label{eq:first_bound_kl}
        \tvnormLigne{p_0 \Pbb_{T|0}- p_{\infty}} \leq \sqrt{2} \exp[-T] \KLLignesqrt{p_0}{p_\infty}  . 
      \end{equation}
      In addition, we have that
      \begin{equation}
        \textstyle{
          \KLLigne{p_0}{p_\infty} = (d/2) \log(2 \uppi) + \int_{\rset^d} \normLigne{x}^2 \rmd p_0(x) - \mathrm{H}(p_0)  ,
          }
        \end{equation}
        where $\mathrm{H}(p_0) = - \int_{\rset^d} \log(p_0(x)) p_0(x) \rmd x$.
        Combining this result and \eqref{eq:first_bound_kl} we get that
        \begin{equation}
          \label{eq:first_bound_kl_clop}
          \textstyle{
            \tvnormLigne{p_0 \Pbb_{T|0}- p_{\infty}} \leq \sqrt{2} \exp[-T] ((d/2) \log(2 \uppi) + \int_{\rset^d} \normLigne{x}^2 \rmd p_0(x) - \mathrm{H}(p_0))^{1/2}   ,
            }
          \end{equation}
          which concludes the first part of the proof.
        \item First, let $\Qbb \in \Pens(\mathcal{C})$ such that
          $\Qbb = p_\infty \Pbb^R_{|0}$, where $\Pbb^R_{|0}$ is the
          disintegration of $\Pbb^R$ w.r.t. $\phi: \ \contspace \to \rset^d$
          given for any $\omega \in \contspace$ by $\phi(\omega) = \omega_T$,
          see \cite{leonard2014some} for instance. Note that for any
          $f \in \rmc(\contspace)$ with $f$ bounded we have
      \begin{align}
        \textstyle{
        \Qbb[f] = \int_{\rset^d} \int_{\contspace} f(\omega) \Pbb^{R}_{|0}(\omega_0, \rmd \omega) \rmd p_\infty(\omega_0)} &= \textstyle{\int_{\rset^d} \int_{\contspace} f(\omega) \Pbb^{R}_{|0}(\omega_0, \rmd \omega) (\rmd p_\infty / \rmd p_T)(\omega_0)\rmd p_T(\omega_0) } \\
&= \textstyle{\int_\contspace f(\omega) (\rmd p_\infty / \rmd p_T)(\omega_0) \rmd \Pbb^R(\omega)}  . 
      \end{align}
      Therefore, we get that for any $\omega \in \contspace$,
      $(\rmd \Qbb / \rmd \Pbb^R)(\omega) = (\rmd p_\infty / \rmd
      p_T)(\omega_0)$.  Let
      $\Rbb = p_\infty \Pbb_{|0}$. Note that for any $t \in \ccint{0,T}$,
      $\Rbb_t = p_\infty$ and that $\Rbb$ is associated with the process
      $\rmd x_t = -x_t \rmd t + \sqrt{2} \rmd w_t$ with
      $\mathcal{L}(x_0) = p_\infty$. In particular, $\Rbb$ satisfies
      \cite[Hypothesis 1.8]{cattiaux2021time}. Using \cite[Theorem
      2.4]{leonard2014some} we have that
      \begin{equation}
        \textstyle{\KLLigne{\Pbb}{\Rbb} = \KLLigne{p_0}{p_\infty} + \int_{\rset^d} \KLLigne{\Pbb_{|0}(x_0)}{\Pbb_{|0}(x_0)} \rmd p_0(x_0) = \KLLigne{p_0}{p_\infty} < +\infty.}
      \end{equation}
      Therefore, we can apply \cite[Theorem 4.9]{cattiaux2021time}. Let
      $u \in \rmc^\infty_c(\rset^d, \rset)$, we have that 
      $(\bfM_t^u(y))_{t \in \ccint{0,T}}$ is a local martingale, where we
      have for any $t \in \ccint{0,T}$
      \begin{equation}
        \textstyle{\bfM_t^u(y) = u(y_t) - u(y_0) - \int_0^t \{ \langle \nabla u(y_s), y_s + 2\nabla \log p_{T-s}(y_s) \rangle + \Delta u(y_s) \}\rmd s ,}
      \end{equation}
      where $\mathcal{L}(y) = \Pbb^R$.  Since $u$ is compactly supported we
      have that
      $\sup_{\omega \in \contspace} \sup_{t \in \ccint{0,T}}
      \absLigne{\bfM_t^u(\omega)} < +\infty$ and therefore
      $(\bfM_t^u(y))_{t \in \ccint{0,T}}$ is a martingale. We now show that
      $(\bfM_t^u(y))_{t \in \ccint{0,T}}$ is a martingale, with
      $\mathcal{L}(y) = \Qbb$. Since
      $\sup_{\omega \in \contspace} \sup_{t \in \ccint{0,T}}
      \absLigne{\bfM_t^u(\omega)} < +\infty$, we have that for any
      $t \in \ \ccint{0,T}$, $\expeLigne{\absLigne{\bfM_t^u}}<+\infty$. Let
      $t,s \in \ccint{0,T}$ with $t > s$ and $g : \contspace \to \rset^d$
      bounded.  We have that
      $\expeLigne{\absLigne{g(\{x_{T-s}\}_{s \in \ccint{0,t}})}^2(\rmd p_\infty / \rmd
        p_T)(x_T)^2} < +\infty$. Hence, we have that
      \begin{equation}
        \expeLigne{(\bfM_t^u(x_{T-\cdot}) - \bfM_s^u(x_{T-\cdot})) g(\{x_{T-s}\}_{s \in \ccint{0,t}})(\rmd p_\infty / \rmd
        p_T)(x_T)} = 0 .
      \end{equation}
      Using this result and that for any $\omega \in \contspace$,
      $(\rmd \Qbb / \rmd \Pbb^R)(\omega) = (\rmd p_\infty / \rmd
      p_T)(\omega_0)$ we get
      \begin{equation}
        \expeLigne{(\bfM_t^u(y) - \bfM_s^u(y)) g(\{y_s\}_{s \in \ccint{0,t}})} = 0 .
      \end{equation}
      Hence, for any $u \in \rmc^2_c(\rset^d, \rset)$,
      $(\bfM_t^u(y))_{t \in \ccint{0,T}}$ is a martingale. In addition,
      $(\bfM_t^u(\bfZ))_{t \in \ccint{0,T}}$ is a martingale using
      \Cref{lemma:strong_solution_backward} and It\^o's lemma, where $\bfZ$ is
      the solution to the SDE in \Cref{lemma:strong_solution_backward}. In
      addition, we have that
      $\mathcal{L}(\bfZ_0) = \mathcal{L}(y_0) = p_\infty$. Using
      \Cref{lemma:control_growth} and the remark following \cite[Hypothesis
      1.8]{cattiaux2021time}, we get that
      $\mathcal{L}(\bfZ) = \mathcal{L}(y) = \Qbb$. We have just shown that the
      time-reversed process with initialisation $p_\infty$ can be obtained as a
      strong solution of an SDE.  Using \Cref{lemma:control_growth} and
      \Cref{lemma:strong_solution_backward}, we have that for any
      $t \in \ccint{0,T}$
      \begin{equation}
        \textstyle{\expeLigne{\int_0^t \normLigne{x_s + 2\nabla \log p_s(x_s)}^2 \rmd s + \int_0^t \normLigne{w_s + 2\nabla \log p_s(w_s)}^2 \rmd s}<+\infty.}
      \end{equation}
      Combining this result and \cite[Lemma S13]{debortoli2021diffusion} we have
      that
      \begin{equation}
        \label{eq:girsanov_uno}
    \textstyle{
    \tvnormLigne{p_\infty \Pbb^R_{T|0} - p_\infty \Qker_0}^2  \leq  (1/2) \int_0^T \expeLigne{\normLigne{b_1(t, (y_s)_{s \in \ccint{0,T}}) - b_2(t, (y_s)_{s \in \ccint{0,T}})}^2} \rmd t  ,}
  \end{equation}
      where for any $t \in \ccint{0,T}$ and $\omega \in \contspace$ we have that
      \begin{equation}
        b_1(t, \omega) = \omega_t + 2\nabla \log p_{T-t}(\omega_t)  , \qquad b_2(t, \omega) = \omega_{t_\gamma} + 2\nabla \log p_{T- t_\gamma}(\omega_{t_\gamma})  ,
      \end{equation}
      where
      $t_\gamma = \sum_{k=0}^{N-1} \mathbb{1}_{\coint{T - t_{k+1}, T -
          t_{k}}}(t) (T - t_{k+1})$. Noting that $(y_t)_{t \in \ccint{0,T}}$
      is distributed according to $\Qbb$ and using that
      $(\rmd \Qbb / \rmd \Pbb^R)(\omega) = (\rmd p_\infty / \rmd
      p_T)(\omega_0)$, \eqref{eq:girsanov_uno} and the Cauchy-Schwarz
      inequality we have
      \begin{align}
        \label{eq:pinsker_girsanov}
        &\textstyle{
          \tvnormLigne{p_\infty \Pbb^R_{T|0} - p_\infty \Qker_0}^2}  \\
        &\textstyle{\leq (1/2) \expeLigne{(\rmd p_\infty / \rmd
          p_T)(x_T)^2}^{1/2} \int_0^T \expeLignesqrt{\normLigne{b_1(t, (x_{T-s})_{s \in \ccint{0,T}}) - b_2(t, (x_{T-s})_{s \in \ccint{0,T}})}^4} \rmd t }  \\
        & \textstyle{\leq (1/2) \expeLigne{(\rmd p_\infty / \rmd
          p_T)(x_T)^2}^{1/2}  }\\
        & \qquad \textstyle{\times \int_0^T \expeLignesqrt{\normLigne{b_1(T-t, (x_{T-s})_{s \in \ccint{0,T}}) - b_2(T-t, (x_{T-s})_{s \in \ccint{0,T}})}^4} \rmd t }  . 
      \end{align}
In addition, we have that for any $t \in \ccint{0,T}$ and $\omega \in \contspace$ we have
\begin{align}
  &\normLigne{b_1(t, \omega)  - b_2(t, \omega)} \\
  & \leq \normLigne{\omega_t - \omega_{t_\gamma}} + 2 \normLigne{\nabla \log p_{T-t}(\omega_t) - \nabla \log p_{T-t_\gamma}(\omega_t)} \\
  & \qquad \qquad + 2 \normLigne{\nabla \log p_{T-t_\gamma}(\omega_t) - \nabla \log p_{T-t_\gamma}(\omega_{t_\gamma})} \\
  &  \textstyle{\leq (1 + 2\sup_{s \in \ccint{0,T}} \sup_{x \in \rset^d}\normLigne{\nabla^2 \log p_s(x)})} \normLigne{\omega_t - \omega_{t_\gamma}} \\
  & \qquad \qquad \textstyle{+ 2 \sup_{s \in \ccint{T-t, T- t_{\gamma}}} \normLigne{\partial_t \nabla \log p_t(\omega_t)}(t - t_\gamma)} \\
&  \textstyle{\leq (1 + 2\Ktt) \normLigne{\omega_t - \omega_{t_\gamma}} + 2 \sup_{s \in \ccint{T-t, T- t_{\gamma}}} \normLigne{\partial_s \nabla \log p_s(\omega_t)} (t - t_\gamma)  }  . 
\end{align}
Note that
\begin{equation}
  \textstyle{
  T - (T - t)_\gamma = T - \sum_{k=0}^{N-1} \mathbb{1}_{\coint{T - t_{k+1}, T -
          t_{k}}}(T - t) (T - t_{k+1}) = \sum_{k=0}^{N-1} \mathbb{1}_{\ocint{t_k,t_{k+1}}}(t) t_{k+1}  . }
\end{equation}
For any $t \in \ccint{0,T}$, denote
$t^\gamma = T - (T - t)_\gamma = \sum_{k=0}^{N-1}
\mathbb{1}_{\ocint{t_k,t_{k+1}}}(t) t_{k+1}$. Therefore, we get that for any
$t \in \ocint{t_k,t_{k+1}}$
\begin{align}
  &\normLigne{b_1(T-t, \omega)  - b_2(T-t, \omega)} \\
  & \qquad \leq \textstyle{(1 + 2 \Ktt) \normLigne{\omega_{T-t} - \omega_{(T-t)_\gamma}} + 2 \sup_{s \in \ccint{t, t^\gamma}} \normLigne{\partial_s \nabla \log p_s(\omega_{T-t})} (t^\gamma - t)} \\
                                                   & \qquad \leq \textstyle{(1 + 2 \Ktt) \normLigne{\omega_{T-t} - \omega_{(T-t)_\gamma}} + 2 \sup_{s \in \ccint{t_k, t_{k+1}}} \normLigne{\partial_s \nabla \log p_s(\omega_{T-t})} \gamma_{k+1}} \\
  & \qquad \leq \textstyle{(1 + 2 \Ktt) \normLigne{\omega_{T-t} - \omega_{(T-t)_\gamma}} + 2 \mathrm{S}_{t_{k}}(\omega_{T-t}) \gamma_{k+1}}  . 
\end{align}
Combining this result and that for any $a,b \geq 0$, $(a+b)^4 \leq 8 a^4 + 8b^4$ we get that for any $t \in \ocint{t_k, t_{k+1}}$
\begin{align}
    \label{eq:almost_there}
    &\expeLigne{\normLigne{b_1(T-t, (x_{T-s})_{s \in \ccint{0,T}})  - b_2(T-t, (x_{T-s})_{s \in \ccint{0,T}})}^4} \\
    & \qquad \qquad \leq \textstyle{8(1 + 2 \Ktt)^4 \expeLigne{\normLigne{x_{t} - x_{t_k}}^4} + 16 \expeLigne{\mathrm{S}_{t_{k}}(x_t)^4} \gamma_{k+1}^4}.
\end{align}
In addition, we have that for any $t \in \ccint{0,T}$, $x_t = \exp[-t]x_0 + w_{(1-\exp[-2t])^{1/2}}$.
Hence, for any $s,t \in \ccint{0,T}$ with $t > s$ we have
\begin{equation}
  \normLigne{x_t - x_s} \leq \exp[-s](\exp[t- s] - 1) \normLigne{x_0} + \normLigne{w_{(1-\exp[-2t])} - w_{(1-\exp[-2s])}}  . 
\end{equation}
Therefore, we have that for any $s,t \in \ccint{0,T}$ with $t > s$
\begin{align}
  \expeLigne{\normLigne{x_t - x_s}^4} &\leq 8\exp[-4s](1 - \exp[-t+s])^4 \expeLigne{\normLigne{x_0}^4} + 8 \expeLigne{\normLigne{w_{(1-\exp[-2t])} - w_{(1-\exp[-2s])}}^4} \\
                                            &\leq 8\exp[-4s](1 - \exp[-t+s])^4 \expeLigne{\normLigne{x_0}^4} + 24 (\exp[-t] - \exp[-s])^2 \\
                                            &\leq 8\exp[-4s](1 - \exp[-t+s])^4 \expeLigne{\normLigne{x_0}^4} + 24 \exp[-2s](1 - \exp[-t+s])^2 \\
  &\leq 8\expeLigne{\normLigne{x_0}^4}\exp[-4s](t-s)^4  + 24 \exp[-2s](t-s)^2  . \label{eq:inter_ou}
\end{align}
In addition, using that that for any $k \in \{0, \dots, N-1\}$ and
$x \in \rset^d$, $\mathtt{S}_{t_k}(x) \leq \Mtt \exp[-\alpha t_k] \normLigne{x}$ we get
that
\begin{equation}
  \expeLigne{\mathtt{S}_{t_k}(x_t)^4} \leq 24 \Mtt^4 \exp[-4\alpha t_k] \{1 + \expeLigne{\normLigne{x_0}^4}\}  . 
\end{equation}

Combining this result, \eqref{eq:almost_there} and \eqref{eq:inter_ou} we get that for any $t \in \ocint{t_k, t_{k+1}}$
\begin{align}
  &\expeLigne{\normLigne{b_1(T-t, (x_{T-s})_{s \in \ccint{0,T}})  - b_2(T-t, (x_{T-s})_{s \in \ccint{0,T}})}^4} \\
  & \qquad \qquad \qquad \leq 64(1 + 2 \Ktt)^4\expeLigne{\normLigne{x_0}^4}\exp[-4t_k]\gamma_{k+1}^4  \\
  &\qquad \qquad \qquad \qquad   + 192 (1 + 2 \Ktt)^4 \exp[-2t_k]\gamma_{k+1}^2 + 384 \Mtt^4 \exp[-4\alpha t_k] \{1 + \expeLigne{\normLigne{x_0}^4}\} \gamma_{k+1}^4  . 
\end{align}
Using this result and that for any $a,b \geq 0$,
$(a+b)^{1/2} \leq a^{1/2} + b^{1/2}$, we have for any $t \in \ocint{t_k,t_{k+1}}$
\begin{align}
  &\expeLignesqrt{\normLigne{b_1(T-t, (x_{T-s})_{s \in \ccint{0,T}})  - b_2(T-t, (x_{T-s})_{s \in \ccint{0,T}})}^4} \\
  & \qquad \qquad \qquad \leq 8(1 + 2 \Ktt)^2\expeLignesqrt{\normLigne{x_0}^4}\exp[-2t_k]\gamma_{k+1}^2  \\
  &\qquad \qquad \qquad  \qquad + 14 (1 + 2 \Ktt)^2 \exp[-t_k]\gamma_{k+1} + 20 \Mtt^2 \exp[-2\alpha t_k] \{1 + \expeLignesqrt{\normLigne{x_0}^4}\} \gamma_{k+1}^2  . \label{eq:almost_almost_there}
\end{align}
We have that for any $\beta >0$,
\begin{equation}
  \textstyle{
 \sum_{k=0}^{N-1} \exp[-\beta t_k] \leq \sum_{k \in \nset} \exp[-\beta
 \gamma_\star k] \leq (1 - \exp[-\beta \gamma_\star])^{-1} \leq 1 + \beta / \gamma_\star  .
 }
\end{equation}
Then using this result, \eqref{eq:almost_almost_there} and \eqref{eq:pinsker_girsanov} we get that
\begin{align}
        &\tvnormLigne{p_\infty \Pbb^R_{T|0} - p_\infty \Qker_0}^2  \textstyle{\leq  \expeLigne{(\rmd p_\infty / \rmd
                                                                        p_T)(x_T)^2}^{1/2}} [4(1 + 2 \Ktt)^2\expeLignesqrt{\normLigne{x_0}^4}(1 +/(2\gamma_\star))(\gamma^\star)^3  \\
  &\qquad \qquad \qquad + 7 (1 + 2 \Ktt)^2 (1 +1/\gamma_\star)(\gamma^\star)^2 + 10 \Mtt^2 \{1 + \expeLignesqrt{\normLigne{x_0}^4}\} (1 +1/(2 \alpha\gamma_\star))  (\gamma^\star)^3]  . 
\end{align}
Therefore, we get that
\begin{align}
        &\tvnormLigne{p_\infty \Pbb^R_{T|0} - p_\infty \Qker_0}  \textstyle{\leq  \expeLigne{(\rmd p_\infty / \rmd
                                                                        p_T)(x_T)^2}^{1/4}} [2(1 + 2 \Ktt)\expeLignesqrtt{\normLigne{x_0}^4}(1 +/(2\gamma_\star)^{1/2})(\gamma^\star)^{3/2}  \\
        &\qquad \qquad \qquad \  + 3 (1 + 2 \Ktt) (1 +1/\gamma_\star^{1/2})\gamma^\star + 4 \Mtt \{1 + \expeLignesqrtt{\normLigne{x_0}^4}\} (1 +1/(2 \alpha\gamma_\star)^{1/2})  (\gamma^\star)^{3/2}] \\
  &\qquad \textstyle{\leq  \expeLigne{(\rmd p_\infty / \rmd
    p_T)(x_T)^2}^{1/4}} [6(1 + 2 \Ktt)(1+\expeLignesqrtt{\normLigne{x_0}^4}) \\
  & \qquad \qquad \qquad + 4 \Mtt \{1 + \expeLignesqrtt{\normLigne{x_0}^4}\} (1 +1/(2 \alpha)^{1/2}) ]((\gamma^\star)^{2}/\gamma_\star)^{1/2} \\
  &\qquad \textstyle{\leq  6 (1+\expeLignesqrtt{\normLigne{x_0}^4}) \expeLigne{(\rmd p_\infty / \rmd
                                                                        p_T)(x_T)^2}^{1/4}} [1 +  \Ktt + \Mtt (1 +1/(2 \alpha)^{1/2}) ]((\gamma^\star)^{2}/\gamma_\star)^{1/2}  ,
\end{align}
which concludes the proof upon using \Cref{item:bound_p_infty}.
    \end{enumerate}    
  \end{proof}

  We now check that the assumption of \Cref{thm:control_diffusion} are satisfied in a Gaussian setting.

\begin{proposition}
  \label{prop:gaussian_case}
  Assume that $p_0 = \mathrm{N}(0, \Sigma)$ and that
  $T \geq 1 +(1/2)[\log^+(\normLigne{\Sigma}) + \log(d+1)]$ then we have that
  for any $t \in \ccint{0,T}$ and $x \in \rset^d$
  \begin{equation}
    \normLigne{\nabla^2 \log p_t(x)} \leq \max(1, \normLigne{\Sigma^{-1}})  , \qquad \normLigne{\partial_t \nabla \log p_t(x)} \leq 2 \exp[-2 t] \max(1, \normLigne{\Sigma^{-1}})^2 \normLigne{\Sigma - \Id}\normLigne{x}  . 
  \end{equation}
  In addition, we have that $\int_{\rset^d} p_\infty(x)^2/p_T(x) \rmd x \leq \sqrt{2}$.
\end{proposition}

\begin{proof}
  Recall that for any $t \in \ccint{0,T}$,
  $x_t = \exp[-t] x_0 + w_{1- \exp[-2t]}$. Therefore, we have that for
  any $t \in \ccint{0,T}$, $p_t = \mathrm{N}(0, \Sigma_t)$ with
  $\Sigma_t = \exp[-2t] \Sigma + (1- \exp[-2t]) \Id$. Hence, we get that for any
  $t \in \ccint{0,T}$ and $x \in \rset^d$,
  $\nabla^2 \log p_t(x) = (\exp[-2t] \Sigma + (1 - \exp[-2t] \Id)^{-1}$. Using
  this result, we have that for any $t \in \ccint{0,T}$ and $x \in \rset^d$,
  $\normLigne{\nabla^2 \log p_t(x)} \leq \max(1, \normLigne{\Sigma^{-1}})$.  Similarly,
  for any $t \in \ccint{0,T}$ and $x \in \rset^d$ we have
  \begin{equation}
    \partial_t \nabla \log p_t(x) = \partial_t \Sigma_t^{-1} x = - \Sigma_t^{-1} (\partial_t \Sigma_t) \Sigma_t^{-1} x  . 
  \end{equation}
  Hence, for any $t \in \ccint{0,T}$ and $x \in \rset^d$ we have
  $\normLigne{\partial_t \nabla \log p_t(x)} \leq 2 \exp[-2 t] \max(1,
  \normLigne{\Sigma^{-1}})^2 \normLigne{\Sigma - \Id}\normLigne{x}$. Finally, we have
  that for any $t \in \ccint{0,T}$ and $x \in \rset^d$
  \begin{equation}
   \langle x, [2 \Id - (\exp[-2t] \Sigma + (1 - \exp[-2t]) \Id)^{-1}]x \rangle \geq (2 - (\exp[-2t] \normLigne{\Sigma^{-1}}^{-1} + (1- \exp[-2t]))^{-1})\normLigne{x}^2  . 
 \end{equation}
 Let $\vareps \in \ocint{0,1/2}$.  For any $t \in \ccint{0,T}$, we have that
 $2 - (\exp[-2t] \normLigne{\Sigma^{-1}}^{-1} + (1- \exp[-2t]))^{-1} \geq 1 -
 \vareps$ if and only if
 $\exp[-2t](1 - \normLigne{\Sigma^{-1}}^{-1}) \leq 1 -
 (1+\vareps)^{-1}$. Using that
 $-\log(1 - (1+\vareps)^{-1}) = \log(1 + \vareps^{-1})$ we have that for any
 $t \geq (1/2) \log(1 + \vareps^{-1})$
 and $x \in \rset^d$
 \begin{equation}
   p_\infty(x)^2/p_t(x) \leq \exp[-\normLigne{x}^2/4](2\uppi)^{-d/2} \det(\Sigma_t)  . 
 \end{equation}
 Combining this result and the fact that
 $\int_{\rset^d} \exp[-\normLigne{x}^2/2(1-\vareps)] \rmd x= (2(1-\vareps)
 \uppi)^{d/2}$, we get that for any 
 $t \geq (1/2) \log(1 +
 \vareps^{-1})$ 
 \begin{equation}
   \textstyle{
     \int_{\rset^d} p_\infty(x)^2/p_t(x) \rmd x \leq \int_{\rset^d} \exp[-\normLigne{x}^2/2(1-\vareps)](2\uppi)^{-d/2} \det(\Sigma_t)^{1/2}  \rmd x \leq (1-\vareps)^{d/2} \det(\Sigma_t)^{1/2} .
     }
   \end{equation}
   Let $\vareps = 1/(2d) \leq 1/2$. Note that
   $T \geq (1/2) \{-\log(\absLigne{\normLigne{\Sigma^{-1}}^{-1}-1}) + \log(1 +
   2d)\}$. Hence, we have that
   \begin{equation}
      \textstyle{
        \int_{\rset^d} p_\infty(x)^2/p_T(x) \rmd x \leq \exp[-\log(1-1/(2d))(d/2)] \det(\Sigma_T)^{1/2} .
        }
   \end{equation}
Since for any $t \in \coint{0,1/2}$, $-\log(1-t)\leq \log(2)t$ we get that
\begin{equation}
  \label{eq:inequality_inter}
    \textstyle{
        \int_{\rset^d} p_\infty(x)^2/p_T(x) \rmd x \leq 2^{1/4} \det(\Sigma_T)^{1/2} .
        }  
\end{equation}
Finally, using that $\Sigma_T = \exp[-2T] \Sigma + (1 - \exp[-2T]) \Id$ we have that
\begin{equation}
  \det(\Sigma_T)^{1/2} \leq (\exp[-2T]\normLigne{\Sigma} + 1 - \exp[-2T])^{d/2} \leq (1 + \exp[-2T] \normLigne{\Sigma})^{d/2} .
\end{equation}
Hence, using that result and that for any $t \geq 0$, $\log(1+t) \leq t$ we have
\begin{equation}
  \det(\Sigma_T)^{1/2} \leq \exp[\exp[-2T]\normLigne{\Sigma}(d/2)] \eqsp . 
\end{equation}
Since,
$T \geq (1/2) \{ \log(\normLigne{\Sigma}) + \log(d) + \log(2) -
\log(\log(2^{1/4})) \}$, we get that $\det(\Sigma_T)^{1/2} \leq 2$, which
concludes the proof upon combining this result and \eqref{eq:inequality_inter}.
\end{proof}

Therefore, we get the following simplified result in the Gaussian setting.

\begin{corollary}
  Assume that $p = \mathrm{N}(0, \Sigma)$, with
  $\normLigne{\Sigma^{-1}} \geq 1$, $\gamma^\star = \gamma_\star= \gamma >0$ and
  $T \geq 1 + (1/2)[\log^+(\normLigne{\Sigma}) + \log(d +1)]$, then we have
  \begin{align}
    \tvnormLigne{\mathcal{L}(X_0) - p_0} &\leq \textstyle{\exp[-T/2] (\log^+(\normLigne{\Sigma^{-1}}) + \normLigne{\Sigma - \Id})^{1/2}} \\
    & \qquad \qquad  \textstyle{+ 48  (1+\normLigne{\Sigma}^{1/2}d^{1/2}) \normLigne{\Sigma^{-1}}^2 [1 +  \normLigne{\Sigma - \Id} ]\sqrt{\gamma}  .} 
  \end{align}  
\end{corollary}

\begin{proof}
  Using \eqref{eq:kl_computation} and \Cref{prop:gaussian_case} we have 
  \begin{align}
    &\tvnormLigne{\mathcal{L}(X_0) - p_0} \leq \textstyle{\exp[-T/2] (-\log(\det(\Sigma)) + \mathrm{Tr}(\Sigma) -d)^{1/2}} \\
    & \qquad \qquad \qquad \qquad \qquad   \textstyle{+ 12  (1+(\int_{\rset^d} \normLigne{x}^4 \rmd p_0(x))^{1/4})  [1 +  \Ktt + 2C ]\sqrt{(\gamma^\star)^{2}/\gamma_\star} } \\
    &\qquad \leq \textstyle{\exp[-T/2] (-\log(\det(\Sigma)) + \mathrm{Tr}(\Sigma) -d)^{1/2}} \\
    & \qquad \qquad  \textstyle{+ 12  (1+(\int_{\rset^d} \normLigne{x}^4 \rmd p_0(x))^{1/4})  [1 +  \normLigne{\Sigma^{-1}} + 2\normLigne{\Sigma^{-1}}^2 \normLigne{\Sigma - \Id} ]\sqrt{(\gamma^\star)^{2}/\gamma_\star}  } \\
    &\qquad \leq \textstyle{\exp[-T/2] (-\log(\det(\Sigma)) + \mathrm{Tr}(\Sigma) -d)^{1/2}} \\
    & \qquad \qquad  \textstyle{+ 12  (1+3^{1/4}\normLigne{\Sigma}^{1/2}d^{1/2})  [1 +  \normLigne{\Sigma^{-1}} + 2\normLigne{\Sigma^{-1}}^2 \normLigne{\Sigma - \Id} ]\sqrt{(\gamma^\star)^{2}/\gamma_\star}  } \\
    &\qquad \leq \textstyle{\exp[-T/2] (-\log(\det(\Sigma)) + \mathrm{Tr}(\Sigma) -d)^{1/2}} \\
    & \qquad \qquad  \textstyle{+ 48  (1+\normLigne{\Sigma}^{1/2}d^{1/2}) \normLigne{\Sigma^{-1}}^2 [1 +  \normLigne{\Sigma - \Id} ]\sqrt{(\gamma^\star)^{2}/\gamma_\star}  } \\
    &\qquad \leq \textstyle{\exp[-T/2] (\log^+(\normLigne{\Sigma^{-1}}) + \normLigne{\Sigma - \Id})^{1/2}} \\
    & \qquad \qquad  \textstyle{+ 48  (1+\normLigne{\Sigma}^{1/2}d^{1/2}) \normLigne{\Sigma^{-1}}^2 [1 +  \normLigne{\Sigma - \Id} ]\sqrt{(\gamma^\star)^{2}/\gamma_\star}  .}             
  \end{align}  
\end{proof}


%%% Local Variables:
%%% mode: latex
%%% TeX-master: "main"
%%% End:
